%%%%%%%%%%%%%%%%%%%%%%%%%%%%%%%%%%%%%%%%%%%%%%%%%%%%%%%%%%%%%%%%%%%%%%%%%%%%%%%%%%
\begin{frame}[fragile]\frametitle{}
\begin{center}
{\Large Introduction}
\end{center}
\end{frame}

%%%%%%%%%%%%%%%%%%%%%%%%%%%%%%%%%%%%%%%%%%%%%%%%%%%%%%%%%%%%%%%%%%%%%%%%%%%%%%%%%%
\begin{frame}[fragile]\frametitle{The World Is a Chessboard}
  \begin{columns}
    \begin{column}[T]{0.55\linewidth}
      \begin{block}{A Strategic Reality}
        We often think life is a series of random events. \textbf{It is not.}
        It is a game, sometimes cooperative, sometimes competitive, but \textit{always strategic}.
      \end{block}
      \vspace{8pt}
      \begin{itemize}
        \item From the vegetable vendor negotiating prices \ldots
        \item \ldots\ to the diplomat averting war.
        \item \textbf{We are all players.}
        \item The question is: are you playing with a strategy?
      \end{itemize}
    \end{column}
    \begin{column}[T]{0.42\linewidth}
      \begin{block}{Who This Is For}
        \begin{itemize}
          \item Professionals \& Managers
          \item Negotiators \& Leaders
          \item Entrepreneurs \& Investors
          \item Curious Minds
          \item \textit{Anyone who interacts with other humans}
        \end{itemize}
      \end{block}
    \end{column}
  \end{columns}
\end{frame}

%%%%%%%%%%%%%%%%%%%%%%%%%%%%%%%%%%%%%%%%%%%%%%%%%%%%%%%%%%%%%%%%%%%%%%%%%%%%%%%%%%
\begin{frame}[fragile]\frametitle{What Is Game Theory?}
  \begin{columns}
    \begin{column}[T]{0.48\linewidth}
      \begin{block}{Definition}
        Game theory is the study of \textbf{strategic interdependence}: situations where your outcome depends not only on what you do, but on what others choose to do.
      \end{block}
      \vspace{8pt}
      \begin{itemize}
        \item \textbf{Players}: Individuals, firms, nations, algorithms
        \item \textbf{Strategies}: Plans of action, sometimes contingent
        \item \textbf{Payoffs}: Outcomes that players value, qualitative or quantitative (money, safety, reputation)
        \item \textbf{Information}: What players know (rarely perfect)
        \item \textbf{Equilibrium}: Stable state where strategies are mutual best responses
      \end{itemize}
    \end{column}
    \begin{column}[T]{0.48\linewidth}
      \begin{block}{Core Insight}
        Success depends on anticipating others' choices and, when possible, shaping them.
      \end{block}
      \vspace{8pt}
      \begin{itemize}
        \item \textbf{Not} about controlling nature or solving physics
        \item \textbf{About} responding to decision-makers who respond
        \item Applies wherever outcomes are shaped by choices
        \item From salary negotiations to nuclear deterrence
        % \item Helps clarify incentives and reveal why outcomes persist
      \end{itemize}
    \end{column}
  \end{columns}
\end{frame}

%%%%%%%%%%%%%%%%%%%%%%%%%%%%%%%%%%%%%%%%%%%%%%%%%%%%%%%%%%%%%%%%%%%%%%%%%%%%%%%%%%
\begin{frame}[fragile]\frametitle{A Brief History}
  \begin{columns}
    \begin{column}[T]{0.5\linewidth}
      \begin{block}{Origins}
        \begin{itemize}
          \item \textbf{1944}, Von Neumann \& Morgenstern publish \textit{Theory of Games and Economic Behaviour}
          \item \textbf{1950}, John Nash introduces the \textit{Nash Equilibrium}
          \item \textbf{1994}, Nash, Harsanyi \& Selten win the Nobel Prize in Economics
          \item \textbf{2005, 2012}, Further Nobel awards for Mechanism Design and Matching Theory
        \end{itemize}
      \end{block}
    \end{column}
    \begin{column}[T]{0.48\linewidth}
      \begin{block}{Key Pioneers}
        \begin{itemize}
          \item \textbf{John von Neumann}, Founder of game theory
          \item \textbf{John Nash}, Equilibrium concept
          \item \textbf{Thomas Schelling}, Focal points \& brinkmanship
          \item \textbf{Kahneman \& Tversky}, Behavioural economics
          \item \textbf{Kautilya (Chanakya)}, Ancient strategic wisdom
        \end{itemize}
      \end{block}
    \end{column}
  \end{columns}
\end{frame}

%%%%%%%%%%%%%%%%%%%%%%%%%%%%%%%%%%%%%%%%%%%%%%%%%%%%%%%%%%%%%%%%%%%%%%%%%%%%%%%%%%
\begin{frame}[fragile]\frametitle{Why Game Theory Matters Today}
  \begin{columns}
    \begin{column}[T]{0.48\linewidth}
      \begin{block}{Traditional Thinking}
        \begin{itemize}
          \item Assumes only you are in control
          \item Optimises in isolation
          \item ``How do I get the best outcome?''
          \item \textit{Like an engineer solving physics}
        \end{itemize}
      \end{block}

    \end{column}
    \begin{column}[T]{0.48\linewidth}
		\begin{block}{Game-Theoretic Thinking}
        \begin{itemize}
          \item Assumes others are strategic too
          \item Optimises in interaction
          \item ``How do I get the best outcome \textbf{given what others will do?}''
          \item \textit{Like a chess player anticipating moves}
        \end{itemize}
      \end{block}
    \end{column}
  \end{columns}
\end{frame}


%%%%%%%%%%%%%%%%%%%%%%%%%%%%%%%%%%%%%%%%%%%%%%%%%%%%%%%%%%%%%%%%%%%%%%%%%%%%%%%%%%
\begin{frame}[fragile]\frametitle{Why Game Theory Matters Today}
  \begin{columns}
    \begin{column}[T]{0.48\linewidth}
      \begin{block}{Where It Applies}
        \begin{itemize}
          \item \textbf{Business}, Pricing, competition, M\&A
          \item \textbf{Politics}, Elections, coalitions, treaties
          \item \textbf{Personal Life}, Negotiation, relationships
          \item \textbf{Technology}, Platform strategy, auctions
          \item \textbf{Society}, Public goods, voting, norms
          \item \textbf{Finance}, Auctions, trading, insurance
        \end{itemize}
      \end{block}
    \end{column}
    \begin{column}[T]{0.48\linewidth}
      \begin{alertblock}{Key Insight}
        You are \textbf{never} playing solitaire (i.e. alone).\\
        You are always playing chess (i.e. against someone).
      \end{alertblock}
    \end{column}
  \end{columns}
\end{frame}

%%%%%%%%%%%%%%%%%%%%%%%%%%%%%%%%%%%%%%%%%%%%%%%%%%%%%%%%%%%%%%%%%%%%%%%%%%%%%%%%%%
\begin{frame}[fragile]\frametitle{How to Use This Playbook}
  \begin{columns}
    \begin{column}[T]{0.48\linewidth}
      \begin{block}{Structure of Each Concept}
        \begin{itemize}
          \item \textbf{Introduction}, What is it?
          \item \textbf{Situation}, When does it apply?
          \item \textbf{Steps}, How to use it?
          \item \textbf{Examples}, Real-world anchors
          \item \textbf{Takeaway}, The bottom line
        \end{itemize}
      \end{block}


    \end{column}
    \begin{column}[T]{0.48\linewidth}
      \begin{block}{Sections Covered}
        \begin{itemize}
          \item Foundational Concepts
          \item Game Types \& Structures
          \item Decision-Making Approaches
          \item Credibility \& Commitment
          \item Strategic Interaction in Relationships
          \item Competitive Advantage
          \item Information \& Uncertainty
          \item Negotiation \& Bargaining
          \item Market Design
          \item Collective Challenges
          \item Prospect Theory
        \end{itemize}
      \end{block}
    \end{column}
  \end{columns}
\end{frame}

%%%%%%%%%%%%%%%%%%%%%%%%%%%%%%%%%%%%%%%%%%%%%%%%%%%%%%%%%%%%%%%%%%%%%%%%%%%%%%%%%%
\begin{frame}[fragile]\frametitle{How to Use This Playbook}
  \begin{columns}
    \begin{column}[T]{0.48\linewidth}
      \begin{block}{A Word of Warning}
        Game Theory is a \textbf{map}, not the territory.
        Real life is messy. People are irrational. Rules change.
        Use these models as lenses, not as algorithms.
      \end{block}
    \end{column}
    \begin{column}[T]{0.48\linewidth}
      \begin{alertblock}{The Infinite Game}
        In the ``Finite Game'' of a football match, you play to \textit{win}.
        In the ``Infinite Game'' of life, you play to \textit{keep playing}.
        Understand the incentives. Respect the players.
      \end{alertblock}
    \end{column}
  \end{columns}
\end{frame}

%%%%%%%%%%%%%%%%%%%%%%%%%%%%%%%%%%%%%%%%%%%%%%%%%%%%%%%%%%%%%%%%%%%%%%%%%%%%%%%%%%
\begin{frame}[fragile]\frametitle{}
\begin{center}
{\Large Foundational Concepts}
\end{center}
\end{frame}

% %%%%%%%%%%%%%%%%%%%%%%%%%%%%%%%%%%%%%%%%%%%%%%%%%%%%%%%%%%%%%%%%%%%%%%%%%%%%%%%%%%
% \begin{frame}[fragile]\frametitle{Understanding Payoff Matrices}
% \begin{columns}
    % \begin{column}[T]{0.5\linewidth}
      % \textbf{What is a Payoff Matrix?}
      % \begin{itemize}
        % \item A table showing outcomes for all strategy combinations
        % \item Rows represent Player 1's outcomes
        % \item Columns represent Player 2's outcomes
        % \item Helps visualize strategic interactions clearly
      % \end{itemize}
      
      % \vspace{6pt}
      % \textbf{Simple Example: Coffee Pricing}
      % \begin{center}
      % \begin{tabular}{|c|c|c|}
        % \hline
        % \textbf{Ramesh $\backslash$ Suresh} & \textbf{High Price} & \textbf{Low Price} \\
        % \hline
        % \textbf{High Price} & (50, 50) & (20, 70) \\
        % \hline
        % \textbf{Low Price} & (70, 20) & (30, 30) \\
        % \hline
      % \end{tabular}
      % \end{center}
      % \small{Payoffs in Rs 1000s of profit}
    % \end{column}
    % \begin{column}[T]{0.5\linewidth}
      % \textbf{How to Read the Matrix:}
      % \begin{itemize}
        % \item \textbf{Identify Players:} Row player (Ramesh) and Column player (Suresh)
        % % \item \textbf{List Strategies:} Each row/column is one strategy option
        % % \item \textbf{Find Outcomes:} Intersection shows what happens when both choose those strategies
        % \item \textbf{Read Payoffs:} First number is always row player's payoff, second is column player's
        % % \item \textbf{Compare:} Look across rows and down columns to find best responses
      % \end{itemize}
    % \end{column}
  % \end{columns}
% \end{frame}

%%%%%%%%%%%%%%%%%%%%%%%%%%%%%%%%%%%%%%%%%%%%%%%%%%%%%%%%%%%%%%%%%%%%%%%%%%%%%%%%%%
\begin{frame}[fragile]\frametitle{Understanding Payoff Matrices}
\begin{columns}
    \begin{column}[T]{0.5\linewidth}
      \textbf{What is a Payoff Matrix?}
      \begin{itemize}
        \item A table showing outcomes for all strategy combinations
        \item Rows represent Player 1's choices (Ramesh)
        \item Columns represent Player 2's choices (Suresh)
        \item Each cell shows the outcome when both make choices
      \end{itemize}
      
      \vspace{6pt}
      \textbf{Example: Coffee Pricing}
      \begin{center}
      \begin{tabular}{|c|c|c|}
        \hline
        \textbf{Ramesh $\backslash$ Suresh} & \textbf{High Price} & \textbf{Low Price} \\
        \hline
        \textbf{High Price} & (50, 50) & (20, 70) \\
        \hline
        \textbf{Low Price} & (70, 20) & (30, 30) \\
        \hline
      \end{tabular}
      \end{center}
      \small{Payoffs in Rs 1000s of profit}
    \end{column}
    \begin{column}[T]{0.5\linewidth}
      \textbf{How to Read the Matrix:}
      \begin{itemize}
        \item First number = Ramesh's profit, second = Suresh's profit
        
        \item \textbf{Example 1:}  
        If both choose \textbf{High Price}, we look at top-left cell:  
        \[
        (50, 50)
        \]
        Both earn Rs 50,000 each
        
        \item \textbf{Example 2:}  
        If Ramesh chooses \textbf{Low Price} and Suresh chooses \textbf{High Price}, we look at:
        \[
        (70, 20)
        \]
        Ramesh earns Rs 70,000, Suresh earns Rs 20,000
        
        \item This shows how one player's decision affects the other
      \end{itemize}
    \end{column}
  \end{columns}
\end{frame}

% %%%%%%%%%%%%%%%%%%%%%%%%%%%%%%%%%%%%%%%%%%%%%%%%%%%%%%%%%%%%%%%%%%%%%%%%%%%%%%%%%%
% \begin{frame}[fragile]\frametitle{Finding Dominant Strategies}
% \begin{columns}
    % \begin{column}[T]{0.5\linewidth}
      % \textbf{What is a Dominant Strategy?}
      % \begin{itemize}
        % \item A strategy that gives you the best payoff \textit{no matter what} the opponent does
        % \item If it exists, it simplifies decision-making dramatically
        % \item You don't need to predict opponent's move
        % \item It's your safest, most robust choice
      % \end{itemize}
      
      % \vspace{6pt}
      % \textbf{Example Matrix:}
      % \begin{center}
      % \begin{tabular}{|c|c|c|}
        % \hline
        % \textbf{You $\backslash$ Opponent} & \textbf{A} & \textbf{B} \\
        % \hline
        % \textbf{Option 1} & (3, 3) & (1, 4) \\
        % \hline
        % \textbf{\textbf{Option 2}} & \textbf{(4, 2)} & \textbf{(2, 3)} \\
        % \hline
      % \end{tabular}
      % \end{center}
    % \end{column}
    % \begin{column}[T]{0.5\linewidth}
      % \textbf{Steps to Find Dominant Strategy:}
      % \begin{itemize}
        % \item \textbf{Pick one of YOUR strategies} (e.g., Option 1)
        % \item \textbf{Compare across ALL opponent strategies:} 
        % \begin{itemize}
          % \item If opponent plays A: You get 3
          % \item If opponent plays B: You get 1
        % \end{itemize}
        % \item \textbf{Repeat for your other strategy} (Option 2):
        % \begin{itemize}
          % \item If opponent plays A: You get 4
          % \item If opponent plays B: You get 2
        % \end{itemize}
        % \item \textbf{Compare:} Option 2 gives 4 vs 3 and 2 vs 1, Option 2 always wins!
        % \item \textbf{Conclusion:} Option 2 is your dominant strategy (bold in matrix)
      % \end{itemize}
    % \end{column}
  % \end{columns}
% \end{frame}

%%%%%%%%%%%%%%%%%%%%%%%%%%%%%%%%%%%%%%%%%%%%%%%%%%%%%%%%%%%%%%%%%%%%%%%%%%%%%%%%%%
\begin{frame}[fragile]\frametitle{Finding Dominant Strategies}
\begin{columns}
    \begin{column}[T]{0.5\linewidth}
      \textbf{What is a Dominant Strategy?}
      \begin{itemize}
        \item A strategy that gives you a better payoff \textit{for every possible move} of your opponent
        \item It works \textbf{no matter what the opponent chooses}
        \item So you don't need to guess or predict anything
      \end{itemize}
      
      \vspace{6pt}
      \textbf{Example Matrix:}
      \begin{center}
      \begin{tabular}{|c|c|c|}
        \hline
        \textbf{You $\backslash$ Opponent} & \textbf{A} & \textbf{B} \\
        \hline
        \textbf{Option 1} & (3, 3) & (1, 4) \\
        \hline
        \textbf{\textbf{Option 2}} & \textbf{(4, 2)} & \textbf{(2, 3)} \\
        \hline
      \end{tabular}
      \end{center}
    \end{column}
    \begin{column}[T]{0.5\linewidth}
      \textbf{How to Think About It:}
      \begin{itemize}
        \item Ask: “If my opponent picks A or B, which option gives me more?”
        
        \item \textbf{Case 1: Opponent chooses A}
        \begin{itemize}
          \item Option 1 gives you 3
          \item Option 2 gives you 4
          \item $\Rightarrow$ Option 2 is better
        \end{itemize}
        
        \item \textbf{Case 2: Opponent chooses B}
        \begin{itemize}
          \item Option 1 gives you 1
          \item Option 2 gives you 2
          \item $\Rightarrow$ Option 2 is again better
        \end{itemize}
        
        \item \textbf{Key Insight:}  
        No matter what opponent does (A or B), Option 2 always gives you a higher payoff
        
        \item \textbf{Conclusion:} Option 2 is a dominant strategy
      \end{itemize}
    \end{column}
  \end{columns}
\end{frame}

% %%%%%%%%%%%%%%%%%%%%%%%%%%%%%%%%%%%%%%%%%%%%%%%%%%%%%%%%%%%%%%%%%%%%%%%%%%%%%%%%%%
% \begin{frame}[fragile]\frametitle{Finding Nash Equilibrium}
% \begin{columns}
    % \begin{column}[T]{0.5\linewidth}
      % \textbf{What is Nash Equilibrium?}
      % \begin{itemize}
        % \item A situation where each player's strategy is the best response to the other's strategy
        % \item No player can improve by changing alone
        % \item The outcome "locks in" even if suboptimal
        % \item May not be the best outcome for both!
      % \end{itemize}
      
      % \vspace{8pt}
      % \textbf{Example Matrix:}
      % \begin{center}
      % \begin{tabular}{|c|c|c|}
        % \hline
        % \textbf{Player 1 $\backslash$ Player 2} & \textbf{Left} & \textbf{Right} \\
        % \hline
        % \textbf{Up} & (5, 5) & (1, 4) \\
        % \hline
        % \textbf{Down} & (4, 1) & \textbf{(2, 2)} \\
        % \hline
      % \end{tabular}
      % \end{center}
    % \end{column}
    % \begin{column}[T]{0.5\linewidth}
      % \textbf{Steps to Find Nash Equilibrium:}
      % \begin{itemize}
        % \item \textbf{For Player 1:} Find best response to each of Player 2's strategies
        % \begin{itemize}
          % \item If P2 plays Left: Up (5) $>$ Down (4)
          % \item If P2 plays Right: Down (2) $>$ Up (1)
        % \end{itemize}
        % \item \textbf{For Player 2:} Find best response to each of Player 1's strategies
        % \begin{itemize}
          % \item If P1 plays Up: Left (5) $>$ Right (4)
          % \item If P1 plays Down: Right (2) $>$ Left (1)
        % \end{itemize}
        % \item \textbf{Find mutual best responses:} (Down, Right) is Nash Equilibrium, both are playing best responses to each other
        % \item \textbf{Note:} Both would prefer (5,5) but can't coordinate!
      % \end{itemize}
    % \end{column}
  % \end{columns}
% \end{frame}

%%%%%%%%%%%%%%%%%%%%%%%%%%%%%%%%%%%%%%%%%%%%%%%%%%%%%%%%%%%%%%%%%%%%%%%%%%%%%%%%%%
\begin{frame}[fragile]\frametitle{Finding Nash Equilibrium: Idea}

\textbf{What is Nash Equilibrium?}
\begin{itemize}
  \item Each player is choosing their \textbf{best response} to the other
  \item No one wants to change their decision \textit{alone}
  \item Think: “Am I doing the best given what the other person is doing?”
\end{itemize}

\vspace{8pt}
\textbf{Example Matrix:}
\begin{center}
\begin{tabular}{|c|c|c|}
  \hline
  \textbf{Player 1 $\backslash$ Player 2} & \textbf{Left} & \textbf{Right} \\
  \hline
  \textbf{Up} & (5, 5) & (1, 4) \\
  \hline
  \textbf{Down} & (4, 1) & (2, 2) \\
  \hline
\end{tabular}
\end{center}

\small{First number = Player 1, second = Player 2}

\end{frame}

%%%%%%%%%%%%%%%%%%%%%%%%%%%%%%%%%%%%%%%%%%%%%%%%%%%%%%%%%%%%%%%%%%%%%%%%%%%%%%%%%%
\begin{frame}[fragile]\frametitle{Step 1: Find Best Responses}

\textbf{For Player 1 (compare numbers in each column):}

\begin{itemize}
  \item If Player 2 plays Left: compare 5 vs 4 $\Rightarrow$ \underline{5 is higher}
  \item If Player 2 plays Right: compare 1 vs 2 $\Rightarrow$ \underline{2 is higher}
\end{itemize}

\vspace{6pt}

\textbf{For Player 2 (compare numbers in each row):}

\begin{itemize}
  \item If Player 1 plays Up: compare 5 vs 4 $\Rightarrow$ \underline{5 is higher}
  \item If Player 1 plays Down: compare 1 vs 2 $\Rightarrow$ \underline{2 is higher}
\end{itemize}

\vspace{10pt}

\textbf{Marking in the matrix:}

\begin{center}
\begin{tabular}{|c|c|c|}
  \hline
  & Left & Right \\
  \hline
  Up & (\underline{5}, \underline{5}) & (1, 4) \\
  \hline
  Down & (4, 1) & (\underline{2}, \underline{2}) \\
  \hline
\end{tabular}
\end{center}

\end{frame}

%%%%%%%%%%%%%%%%%%%%%%%%%%%%%%%%%%%%%%%%%%%%%%%%%%%%%%%%%%%%%%%%%%%%%%%%%%%%%%%%%%
\begin{frame}[fragile]\frametitle{Step 2: Find Nash Equilibrium}

\textbf{Key Idea:}
\begin{itemize}
  \item Nash equilibrium occurs where \textbf{both players are choosing best responses}
  \item That means: both numbers in the cell are “selected” (underlined)
\end{itemize}

\vspace{6pt}

\begin{center}
\begin{tabular}{|c|c|c|}
  \hline
  & Left & Right \\
  \hline
  Up & (\underline{5}, \underline{5}) & (1, 4) \\
  \hline
  Down & (4, 1) & (\underline{2}, \underline{2}) \\
  \hline
\end{tabular}
\end{center}

\vspace{6pt}

\textbf{Observation:}
\begin{itemize}
  \item Two cells have both values underlined:
  \begin{itemize}
    \item (Up, Left) $\Rightarrow$ (5,5)
    \item (Down, Right) $\Rightarrow$ (2,2)
  \end{itemize}
\end{itemize}

\vspace{6pt}

\textbf{Conclusion:}
\begin{itemize}
  \item This game has \textbf{two Nash equilibria}
\end{itemize}

\end{frame}

%%%%%%%%%%%%%%%%%%%%%%%%%%%%%%%%%%%%%%%%%%%%%%%%%%%%%%%%%%%%%%%%%%%%%%%%%%%%%%%%%%
\begin{frame}[fragile]\frametitle{Is (5,5) a Nash Equilibrium?}

\textbf{Short Answer: Yes, it \textit{is} a Nash Equilibrium}

\vspace{6pt}

\textbf{Why?}
\begin{itemize}
  \item At (Up, Left):
  \begin{itemize}
    \item Player 1 gets 5 — switching to Down gives 4 (worse)
    \item Player 2 gets 5 — switching to Right gives 4 (worse)
  \end{itemize}
  \item So neither player wants to deviate
\end{itemize}

\vspace{8pt}

\textbf{Then why focus on (2,2)?}
\begin{itemize}
  \item (2,2) is also a Nash equilibrium
  \item But it is \textbf{worse for both players}
  \item This shows:
  \begin{itemize}
    \item Nash equilibrium is about \textbf{stability}, not \textbf{optimality}
  \end{itemize}
\end{itemize}

\vspace{6pt}

\textbf{Key Insight:}
\begin{itemize}
  \item Players may get stuck in a worse outcome even when a better one exists
\end{itemize}

\end{frame}

%%%%%%%%%%%%%%%%%%%%%%%%%%%%%%%%%%%%%%%%%%%%%%%%%%%%%%%%%%%%%%%%%%%%%%%%%%%%%%%%%%
\begin{frame}[fragile]\frametitle{Which Nash Equilibrium Will Be Chosen?}

\textbf{Recall: We found two Nash equilibria}

\begin{center}
\begin{tabular}{|c|c|c|}
  \hline
  & Left & Right \\
  \hline
  Up & (\underline{5}, \underline{5}) & (1, 4) \\
  \hline
  Down & (4, 1) & (\underline{2}, \underline{2}) \\
  \hline
\end{tabular}
\end{center}

\vspace{6pt}

\textbf{Question: Which outcome will players choose?}

\vspace{6pt}

\begin{itemize}
  \item Both (5,5) and (2,2) are stable (no one wants to deviate alone)
  
  \item \textbf{(5,5) is better for both}  
  \begin{itemize}
    \item Higher payoff for Player 1 and Player 2
  \end{itemize}
  
  \item But there is a problem: \textbf{coordination}
  \begin{itemize}
    \item To reach (5,5), both must choose (Up, Left)
    \item If one player makes a mistake, payoff drops sharply
  \end{itemize}
  
  \item \textbf{(2,2) is safer}
  \begin{itemize}
    \item Even if unsure about the other player, risk is lower
  \end{itemize}
\end{itemize}

\vspace{6pt}

\textbf{Key Insight:}
\begin{itemize}
  \item Some equilibria are \textbf{better}, others are \textbf{safer}
  \item Real-world outcomes depend on \textbf{trust, expectations, and coordination}
\end{itemize}

\end{frame}

%%%%%%%%%%%%%%%%%%%%%%%%%%%%%%%%%%%%%%%%%%%%%%%%%%%%%%%%%%%%%%%%%%%%%%%%%%%%%%%%%%
\begin{frame}[fragile]\frametitle{Coordination Problem: Risk vs Reward}

\textbf{Think like a player:}

\begin{itemize}
  \item “I would like (5,5)… but what if the other player doesn’t cooperate?”
\end{itemize}

\vspace{6pt}

\textbf{Compare the risk:}

\begin{itemize}
  \item If you aim for (5,5) but the other player deviates:
  \begin{itemize}
    \item You may end up with \textbf{1} (very low payoff)
  \end{itemize}
  
  \item If you aim for (2,2) and the other deviates:
  \begin{itemize}
    \item Your payoff does not fall as drastically
  \end{itemize}
\end{itemize}

\vspace{8pt}

\textbf{Interpretation:}
\begin{itemize}
  \item (5,5) = \textbf{High reward, high risk}
  \item (2,2) = \textbf{Lower reward, safer choice}
\end{itemize}

\vspace{8pt}

\textbf{Real-world analogy:}
\begin{itemize}
  \item Two firms setting high prices (cooperate) vs undercutting each other
  \item Trust leads to better outcomes, but fear leads to safer ones
\end{itemize}

\end{frame}

%%%%%%%%%%%%%%%%%%%%%%%%%%%%%%%%%%%%%%%%%%%%%%%%%%%%%%%%%%%%%%%%%%%%%%%%%%%%%%%%%%
\begin{frame}[fragile]\frametitle{How to Reach the Better Equilibrium (5,5)?}

\textbf{Problem:}
\begin{itemize}
  \item (5,5) is better for both players
  \item But it requires \textbf{coordination and trust}
\end{itemize}

\vspace{8pt}

\textbf{What can help players coordinate?}

\begin{itemize}
  \item \textbf{Communication}
  \begin{itemize}
    \item Players can agree beforehand: “Let’s both choose (Up, Left)”
  \end{itemize}
  
  \item \textbf{Repeated Interaction}
  \begin{itemize}
    \item If the game is played many times, cooperation can be rewarded
    \item Deviating once may lead to punishment later
  \end{itemize}
  
  \item \textbf{Trust / Reputation}
  \begin{itemize}
    \item If a player is known to cooperate, others are more likely to do the same
  \end{itemize}
  
  \item \textbf{Focal Points (Obvious Choices)}
  \begin{itemize}
    \item Players may naturally gravitate to “prominent” outcomes like (5,5)
  \end{itemize}
\end{itemize}

\vspace{6pt}

\textbf{Key Idea:}
\begin{itemize}
  \item Better outcomes often require \textbf{coordination mechanisms}, not just rational thinking
\end{itemize}

\end{frame}

%%%%%%%%%%%%%%%%%%%%%%%%%%%%%%%%%%%%%%%%%%%%%%%%%%%%%%%%%%%%%%%%%%%%%%%%%%%%%%%%%%
\begin{frame}[fragile]\frametitle{What If Coordination Fails?}

\textbf{Without coordination:}

\begin{itemize}
  \item Players may avoid risk and choose safer strategies
  \item This can lead to the equilibrium:
  \[
  (2,2)
  \]
\end{itemize}

\vspace{8pt}

\textbf{Why does this happen?}

\begin{itemize}
  \item Fear of mismatch:
  \begin{itemize}
    \item “What if I choose Up but the other chooses Right?”
  \end{itemize}
  
  \item Lack of trust:
  \begin{itemize}
    \item No guarantee the other player will cooperate
  \end{itemize}
  
  \item No communication:
  \begin{itemize}
    \item Players cannot coordinate their choices
  \end{itemize}
\end{itemize}

\vspace{8pt}

\textbf{Final Insight:}
\begin{itemize}
  \item Rational players can end up in \textbf{worse outcomes}
  \item The issue is not intelligence, but \textbf{lack of coordination}
\end{itemize}

\end{frame}

%%%%%%%%%%%%%%%%%%%%%%%%%%%%%%%%%%%%%%%%%%%%%%%%%%%%%%%%%%%%%%%%%%%%%%%%%%%%%%%%%%
\begin{frame}[fragile]\frametitle{Real-World Example: Traffic Coordination}

\textbf{Imagine: Choosing which side of the road to drive on}

\begin{center}
\begin{tabular}{|c|c|c|}
  \hline
  \textbf{Driver 1 $\backslash$ Driver 2} & \textbf{Left} & \textbf{Right} \\
  \hline
  \textbf{Left} & (5, 5) & (0, 0) \\
  \hline
  \textbf{Right} & (0, 0) & (5, 5) \\
  \hline
\end{tabular}
\end{center}

\vspace{6pt}

\begin{itemize}
  \item If both choose the \textbf{same side}, traffic flows smoothly → high payoff (5,5)
  
  \item If they choose \textbf{different sides}, there is chaos (or accidents!) → payoff (0,0)
\end{itemize}

\vspace{8pt}

\textbf{Key Observations:}
\begin{itemize}
  \item There are \textbf{two Nash equilibria}:
  \begin{itemize}
    \item (Left, Left)
    \item (Right, Right)
  \end{itemize}
  
  \item Both are equally good—but players must \textbf{coordinate}
\end{itemize}

\end{frame}

%%%%%%%%%%%%%%%%%%%%%%%%%%%%%%%%%%%%%%%%%%%%%%%%%%%%%%%%%%%%%%%%%%%%%%%%%%%%%%%%%%
\begin{frame}[fragile]\frametitle{Why Do We Successfully Coordinate?}

\textbf{Question: Why don’t drivers crash every day?}

\vspace{6pt}

\begin{itemize}
  \item \textbf{Social conventions}
  \begin{itemize}
    \item In India, everyone follows “drive on the left”
  \end{itemize}
  
  \item \textbf{Clear rules and laws}
  \begin{itemize}
    \item Government enforces one equilibrium
  \end{itemize}
  
  \item \textbf{Common expectations}
  \begin{itemize}
    \item You expect others to follow the same rule
  \end{itemize}
\end{itemize}

\vspace{8pt}

\textbf{Connection to Game Theory:}
\begin{itemize}
  \item Real-world systems often \textbf{select one equilibrium}
  \item Once established, everyone follows it automatically
\end{itemize}

\vspace{8pt}

\textbf{Final Insight:}
\begin{itemize}
  \item Good outcomes are not just about incentives
  \item They depend on \textbf{rules, norms, and coordination mechanisms}
\end{itemize}

\end{frame}


%%%%%%%%%%%%%%%%%%%%%%%%%%%%%%%%%%%%%%%%%%%%%%%%%%%%%%%%%%%%%%%%%%%%%%%%%%%%%%%%%%
\begin{frame}[fragile]\frametitle{}
\begin{center}
{\Large Popular Strategies}
\end{center}
\end{frame}

%%%%%%%%%%%%%%%%%%%%%%%%%%%%%%%%%%%%%%%%%%%%%%%%%%%%%%%%%%%%%%%%%%%%%%%%%%%%%%%%%%
\begin{frame}[fragile]\frametitle{Nash Equilibrium: The Stable State}
\begin{columns}
    \begin{column}[T]{0.5\linewidth}
      \begin{itemize}
        \item \textbf{Introduction}: A state where no player can improve their outcome by changing their strategy alone. It is the point at which the game \textit{self-locks}, often into a suboptimal outcome for all.
        \item \textbf{Situation}: Arises in any repeated interaction where actors respond to each other, arms races, price wars, social norms, congestion. Whenever you are stuck in a bad pattern everyone perpetuates.
      \end{itemize}
    \end{column}
    \begin{column}[T]{0.5\linewidth}
      \begin{itemize}
        \item \textbf{Steps}:
        \begin{itemize}
          \item List all players and their possible strategies
          \item Map payoffs for each combination
          \item For each player, find the best response to every opponent strategy
          \item Identify where best responses align, that is the equilibrium
          \item Ask: is it efficient? If not, what changes the game?
        \end{itemize}\item \textbf{Examples}: Everyone stands at a cricket match so no one sees better; arms races leaving all nations poorer; airlines locked in price wars destroying profitability.
      \end{itemize}
    \end{column}
  \end{columns}
\end{frame}

%%%%%%%%%%%%%%%%%%%%%%%%%%%%%%%%%%%%%%%%%%%%%%%%%%%%%%%%%%%%%%%%%%%%%%%%%%%%%%%%%%
\begin{frame}[fragile]\frametitle{Nash Equilibrium: Ramesh-Suresh Coffee Shop Example}
\begin{columns}
    \begin{column}[T]{0.5\linewidth}
      \textbf{The Scenario:}
      \begin{itemize}
        \item Two coffee shops across the street from each other
        \item Originally both charged Rs 100 per cup, profits were stable
        \item Each can independently choose to charge Rs 100 or Rs 70
        \item If one cuts price while other maintains, the cutter gains customers
        \item Both are stuck at Rs 70 now : neither can improve alone
      \end{itemize}
      
      \vspace{2pt}
      \textbf{Payoff Matrix (Daily Profit):}
      \begin{center}
      \begin{tabular}{|c|c|c|}
        \hline
        \textbf{Ramesh $\backslash$ Suresh} & \textbf{Rs 100} & \textbf{Rs 70} \\
        \hline
        \textbf{Rs 100} & (5k, 5k) & (1k, 7k) \\
        \hline
        \textbf{Rs 70} & (7k, 1k) & \textbf{(3k, 3k)} \\
        \hline
      \end{tabular}
      \end{center}
    \end{column}
    \begin{column}[T]{0.5\linewidth}
      \textbf{Analysis Steps:}
      \begin{itemize}
        \item \textbf{Check Best Responses:}
        \begin{itemize}
          \item If Suresh charges Rs 100, Ramesh's best = Rs 70 (7k $>$ 5k)
          \item If Suresh charges Rs 70, Ramesh's best = Rs 70 (3k $>$ 1k)
          \item Same logic for Suresh
        \end{itemize}
        \item \textbf{Nash Equilibrium:} Both charge Rs 70 (bold in matrix)
        \item \textbf{The Trap:} Each earns only Rs 3000 instead of Rs 5000
        \item \textbf{Key Insight:} No player can improve by changing alone : the equilibrium is stable but suboptimal
      \end{itemize}
    \end{column}
  \end{columns}
\end{frame}

%%%%%%%%%%%%%%%%%%%%%%%%%%%%%%%%%%%%%%%%%%%%%%%%%%%%%%%%%%%%%%%%%%%%%%%%%%%%%%%%%%
\begin{frame}[fragile]\frametitle{Dominant Strategy: The Safe Move}
\begin{columns}
    \begin{column}[T]{0.5\linewidth}
      \begin{itemize}
        \item \textbf{Introduction}: A dominant strategy delivers the highest payoff for a player \textit{regardless of what others do}. It is the ``no-brainer'', you need not read the opponent's mind when you have one.
        \item \textbf{Situation}: Look for dominant strategies before trying to anticipate opponents. Useful in competitive exams, auctions, market pricing, and any situation where one choice outperforms all others unconditionally.
      \end{itemize}
    \end{column}
    \begin{column}[T]{0.5\linewidth}
      \begin{itemize}
        \item \textbf{Steps}:
        \begin{itemize}
          \item Draw a payoff matrix
          \item For each of your strategies, find the payoffs against every opponent move
          \item Check if one row always beats the others
          \item If yes, that is your dominant strategy, play it
          \item If not, identify your opponent's dominant strategy and best-respond to it
        \end{itemize}\item \textbf{Examples}: Amazon's relentless cost leadership; index fund investing over stock-picking; vaccination during a pandemic regardless of what neighbours do.
      \end{itemize}
    \end{column}
  \end{columns}
\end{frame}

%%%%%%%%%%%%%%%%%%%%%%%%%%%%%%%%%%%%%%%%%%%%%%%%%%%%%%%%%%%%%%%%%%%%%%%%%%%%%%%%%%
\begin{frame}[fragile]\frametitle{Dominant Strategy: Ramesh-Suresh Wi-Fi Example}
\begin{columns}
    \begin{column}[T]{0.55\linewidth}
      \textbf{The Scenario:}
      \begin{itemize}
        \item Ramesh considers adding high-speed Wi-Fi to his coffee shop
        \item If Suresh doesn't install Wi-Fi, Ramesh attracts students \& freelancers
        \item If Suresh does install Wi-Fi, Ramesh needs it to stay competitive
        \item Either way, installing Wi-Fi makes Ramesh better off
        \item This is a dominant strategy : works regardless of opponent's choice
      \end{itemize}
      
      \vspace{2pt}
      \textbf{Payoff Matrix (Weekly Additions):}
      \begin{center}
      \begin{tabular}{|c|c|c|}
        \hline
        \textbf{Ramesh $\backslash$ Suresh} & \textbf{No Wi-Fi} & \textbf{Install Wi-Fi} \\
        \hline
        \textbf{No Wi-Fi} & (50, 50) & (20, 90) \\
        \hline
        \textbf{\textbf{Install Wi-Fi}} & \textbf{(100, 30)} & \textbf{(60, 60)} \\
        \hline
      \end{tabular}
      \end{center}
    \end{column}
    \begin{column}[T]{0.45\linewidth}
      \textbf{Analysis Steps:}
      \begin{itemize}
        \item \textbf{Test Ramesh's Options:}
        \begin{itemize}
          \item If Suresh has no Wi-Fi: Install (100) $>$ Don't (50)
          \item If Suresh has Wi-Fi: Install (60) $>$ Don't (20)
        \end{itemize}
        \item \textbf{Dominant Strategy:} Install Wi-Fi (bold row) : always better regardless of Suresh's choice
        \item \textbf{Same Logic for Suresh:} Installing Wi-Fi dominates for him too
        \item \textbf{Key Insight:} No mind-reading required. The strategy is robust : it's insurance, not cleverness
      \end{itemize}
    \end{column}
  \end{columns}
\end{frame}

%%%%%%%%%%%%%%%%%%%%%%%%%%%%%%%%%%%%%%%%%%%%%%%%%%%%%%%%%%%%%%%%%%%%%%%%%%%%%%%%%%
\begin{frame}[fragile]\frametitle{Zero-Sum Games: Winner-Takes-All}
\begin{columns}
    \begin{column}[T]{0.5\linewidth}
      \begin{itemize}
        \item \textbf{Introduction}: In a zero-sum game the total available benefit is fixed, one player's gain is another's exact loss. These are games of pure conflict where cooperation is mathematically impossible.
        \item \textbf{Situation}: Applicable where resources are genuinely scarce and non-expandable: elections, IIT seats, market share in a stagnant industry, athletic competitions, inheritance division.
      \end{itemize}
    \end{column}
    \begin{column}[T]{0.5\linewidth}
      \begin{itemize}
        \item \textbf{Steps}:
        \begin{itemize}
          \item Verify the game is truly zero-sum before treating it as one
          \item Identify your relative strengths vs.\ opponents
          \item Use minimax strategy, minimise maximum possible loss
          \item Look for dominant positions to exploit
          \item Consider whether the game can be redesigned as positive-sum
        \end{itemize}\item \textbf{Examples}: Poker tables; election seat counts; dividing a fixed estate; ranking systems where one rise means another's fall.
      \end{itemize}
    \end{column}
  \end{columns}
\end{frame}


%%%%%%%%%%%%%%%%%%%%%%%%%%%%%%%%%%%%%%%%%%%%%%%%%%%%%%%%%%%%%%%%%%%%%%%%%%%%%%%%%%
\begin{frame}[fragile]\frametitle{Zero-Sum Games: Ramesh-Suresh Customer War}
\begin{columns}
    \begin{column}[T]{0.58\linewidth}
      \textbf{The Scenario:}
      \begin{itemize}
        \item Both locked in war over same 50 factory workers who cross street every morning
        \item Fixed customer pool, when one gains a customer, other loses one
        \item Discounts, louder music, bigger signboards
        \item Every move designed to steal, not to create
        \item Victory for one = defeat for other
        \item Classic zero-sum: customer pool is fixed
      \end{itemize}
      
      \vspace{2pt}
      \textbf{Payoff Matrix (Daily Customers):}
      \begin{center}
      \begin{tabular}{|c|c|c|}
        \hline
        \textbf{Ramesh $\backslash$ Suresh} & \textbf{Normal} & \textbf{Aggressive} \\
        \hline
        \textbf{Normal} & (25, 25) & (15, 35) \\
        \hline
        \textbf{Aggressive} & (35, 15) & (25, 25) \\
        \hline
      \end{tabular}
      \end{center}
    \end{column}
    \begin{column}[T]{0.42\linewidth}
      \textbf{Analysis Steps:}
      \begin{itemize}
        \item \textbf{Check Total:} Each cell sums to 50 customers (25+25, 15+35, etc.)
        \item \textbf{Zero-Sum Property:} One's gain = Other's exact loss
        \item \textbf{No Win-Win:} Cannot both gain customers simultaneously
        \item \textbf{Pure Conflict:} Cooperation doesn't expand the pie
        \item \textbf{Key Insight:} Verify game is truly zero-sum before treating it as one, many situations that seem zero-sum are actually positive-sum
      \end{itemize}
    \end{column}
  \end{columns}
\end{frame}


%%%%%%%%%%%%%%%%%%%%%%%%%%%%%%%%%%%%%%%%%%%%%%%%%%%%%%%%%%%%%%%%%%%%%%%%%%%%%%%%%%
\begin{frame}[fragile]\frametitle{Positive-Sum Games: Creating Mutual Benefit}
\begin{columns}
    \begin{column}[T]{0.5\linewidth}
      \begin{itemize}
        \item \textbf{Introduction}: Positive-sum games allow the total gains to grow, ``growing the pie'' rather than fighting over slices. Trade, collaboration, and innovation are classic examples where mutual gain is possible.
        \item \textbf{Situation}: Most business partnerships, trade agreements, marriages, and open-source projects. When both parties gain from interaction, the strategic goal shifts from defeating to creating and coordinating.
      \end{itemize}
    \end{column}
    \begin{column}[T]{0.5\linewidth}
      \begin{itemize}
        \item \textbf{Steps}:
        \begin{itemize}
          \item Map potential joint gains from cooperation
          \item Identify each party's comparative advantage
          \item Design a split of gains that both prefer over non-cooperation
          \item Build trust mechanisms to prevent defection
          \item Iterate, reputation and repeat interaction sustain positive-sum outcomes
        \end{itemize}\item \textbf{Examples}: Bangalore's tech ecosystem where new firms benefit incumbents; open-source software; international trade agreements; ridesharing platforms benefiting both drivers and riders.
      \end{itemize}
    \end{column}
  \end{columns}
\end{frame}


%%%%%%%%%%%%%%%%%%%%%%%%%%%%%%%%%%%%%%%%%%%%%%%%%%%%%%%%%%%%%%%%%%%%%%%%%%%%%%%%%%
\begin{frame}[fragile]\frametitle{Positive-Sum: Ramesh-Suresh Grow Together}
\begin{columns}
    \begin{column}[T]{0.6\linewidth}
      \textbf{The Scenario:}
      \begin{itemize}
        \item Both improve coffee, seating, vibe
        \item Street becomes destination
        \item Office workers detour on purpose, weekend crowds arrive
        \item Better beans from suppliers, local bakery partnerships
        \item Each one's improvements lift expectations for both
        \item Total customers increase instead of staying fixed
      \end{itemize}
      
      \vspace{6pt}
      \textbf{Payoff Matrix (Daily Customers):}
      \begin{center}
      \begin{tabular}{|c|c|c|}
        \hline
        \textbf{Ramesh $\backslash$ Suresh} & \textbf{Low Quality} & \textbf{High Quality} \\
        \hline
        \textbf{Low Quality} & (25, 25) & (30, 40) \\
        \hline
        \textbf{High Quality} & (40, 30) & \textbf{(55, 55)} \\
        \hline
      \end{tabular}
      \end{center}
    \end{column}
    \begin{column}[T]{0.38\linewidth}
      \textbf{Analysis Steps:}
      \begin{itemize}
        \item \textbf{Check Totals:}
        \begin{itemize}
          \item Both Low: 50 total
          \item Mixed: 70 total
          \item Both High: 110 total!
        \end{itemize}
        \item \textbf{Positive-Sum Property:} Pie grows with cooperation
        \item \textbf{Best Outcome:} Both High Quality (55, 55)
        \item \textbf{Key Insight:} Real question isn't how to beat rival, it's how to grow the total market
        \item \textbf{Biggest wins:} Come from making location a destination
      \end{itemize}
    \end{column}
  \end{columns}
\end{frame}


% %%%%%%%%%%%%%%%%%%%%%%%%%%%%%%%%%%%%%%%%%%%%%%%%%%%%%%%%%%%%%%%%%%%%%%%%%%%%%%%%%%
% \begin{frame}[fragile]\frametitle{}
% \begin{center}
% {\Large Game Types and Structures}
% \end{center}
% \end{frame}

%%%%%%%%%%%%%%%%%%%%%%%%%%%%%%%%%%%%%%%%%%%%%%%%%%%%%%%%%%%%%%%%%%%%%%%%%%%%%%%%%%
\begin{frame}[fragile]\frametitle{Prisoner's Dilemma: The Trust Problem}
\begin{columns}
    \begin{column}[T]{0.5\linewidth}
      \begin{itemize}
        \item \textbf{Introduction}: The most famous paradox in game theory, individual rationality (self-protection) leads to collective irrationality (both suffer). Mutual cooperation would be best, but each defects, fearing betrayal.
        \item \textbf{Situation}: Any situation where independent self-interest produces a worse collective outcome: price wars, arms races, doping in sports, environmental pollution, negative political campaigning.
      \end{itemize}
    \end{column}
    \begin{column}[T]{0.5\linewidth}
      \begin{itemize}
        \item \textbf{Steps}:
        \begin{itemize}
          \item Recognise the dilemma structure in your situation
          \item Assess the shadow of the future, is there a next interaction?
          \item Build communication and trust mechanisms
          \item Introduce contracts or penalties for defection
          \item Use tit-for-tat strategies in repeated play
        \end{itemize}\item \textbf{Examples}: Two sweet shops undercutting each other to thin margins; competing nations in arms races; athletes using performance-enhancing drugs; public littering.
      \end{itemize}
    \end{column}
  \end{columns}
\end{frame}

%%%%%%%%%%%%%%%%%%%%%%%%%%%%%%%%%%%%%%%%%%%%%%%%%%%%%%%%%%%%%%%%%%%%%%%%%%%%%%%%%%
\begin{frame}[fragile]\frametitle{Prisoner's Dilemma: Ramesh-Suresh Price War}
\begin{columns}
    \begin{column}[T]{0.6\linewidth}
      \textbf{The Scenario:}
      \begin{itemize}
        \item Both shops initially charge fair prices and maintain quality
        \item Each can choose: \textit{Cooperate} (fair pricing) or \textit{Defect} (cut prices aggressively)
        \item If both cooperate: stable profits, good margins
        \item If one defects while other cooperates: defector gains market share
        \item If both defect: price war destroys profitability for both
        \item Fear of betrayal drives both to defect
      \end{itemize}
      
      \vspace{2pt}
      \textbf{Payoff Matrix (Monthly Profit in Rs):}
      \begin{center}
      \begin{tabular}{|c|c|c|}
        \hline
        \textbf{Ramesh $\backslash$ Suresh} & \textbf{Cooperate} & \textbf{Defect} \\
        \hline
        \textbf{Cooperate} & (50k, 5k0) & (20k, 70k) \\
        \hline
        \textbf{Defect} & (70k, 20k) & \textbf{(30k, 30k)} \\
        \hline
      \end{tabular}
      \end{center}
    \end{column}
    \begin{column}[T]{0.38\linewidth}
      \textbf{Analysis Steps:}
      \begin{itemize}
        \item \textbf{Check Incentives:}
        \begin{itemize}
          \item If Suresh cooperates: Defect (70K) $>$ Cooperate (50K)
          \item If Suresh defects: Defect (30K) $>$ Cooperate (20K)
        \end{itemize}
        \item \textbf{Dominant Strategy:} Defect for both players
        \item \textbf{Nash Equilibrium:} (Defect, Defect) : bold in matrix
        \item \textbf{The Paradox:} Both earn 30K instead of 50K
        \item \textbf{Escape Routes:} Repeated interaction, trust-building, contracts, tit-for-tat
      \end{itemize}
    \end{column}
  \end{columns}
\end{frame}

%%%%%%%%%%%%%%%%%%%%%%%%%%%%%%%%%%%%%%%%%%%%%%%%%%%%%%%%%%%%%%%%%%%%%%%%%%%%%%%%%%
\begin{frame}[fragile]\frametitle{Chicken Game: The Dangerous Standoff}
\begin{columns}
    \begin{column}[T]{0.5\linewidth}
      \begin{itemize}
        \item \textbf{Introduction}: A test of nerves where the worst outcome is mutual destruction. Unlike the Prisoner's Dilemma, someone \textit{must} yield. The winner is the one who convincingly cannot back down.
        \item \textbf{Situation}: High-stakes confrontations where one party must blink: geopolitical crises, union vs.\ management standoffs, legislative gridlock, corporate takeover battles, sovereign debt negotiations.
      \end{itemize}
    \end{column}
    \begin{column}[T]{0.5\linewidth}
      \begin{itemize}
        \item \textbf{Steps}:
        \begin{itemize}
          \item Assess the cost of mutual destruction to both sides
          \item Signal commitment credibly, make retreat costly or visible
          \item Look for face-saving offramps for the opponent
          \item Avoid symmetrical escalation
          \item Know your red lines before the confrontation begins
        \end{itemize}\item \textbf{Examples}: Cuban Missile Crisis; union strikes vs.\ management lockouts; government debt-ceiling standoffs; two buses on a narrow road.
      \end{itemize}
    \end{column}
  \end{columns}
\end{frame}


%%%%%%%%%%%%%%%%%%%%%%%%%%%%%%%%%%%%%%%%%%%%%%%%%%%%%%%%%%%%%%%%%%%%%%%%%%%%%%%%%%
\begin{frame}[fragile]\frametitle{Stag Hunt: Coordination Through Trust}
\begin{columns}
    \begin{column}[T]{0.5\linewidth}
      \begin{itemize}
        \item \textbf{Introduction}: A coordination game with two equilibria, a superior cooperative outcome (the Stag) and an inferior but safe solo outcome (the Hare). The challenge is not conflict but \textit{assurance}: knowing others will commit.
        \item \textbf{Situation}: Climate accords, new technology standards, currency adoption, startup equity decisions, team sports defence, any endeavour where joint success requires full participation.
      \end{itemize}
    \end{column}
    \begin{column}[T]{0.5\linewidth}
      \begin{itemize}
        \item \textbf{Steps}:
        \begin{itemize}
          \item Communicate the payoff difference between cooperation and going solo
          \item Build pre-commitment mechanisms
          \item Start with a small, credible coalition
          \item Provide insurance against defection
          \item Create repeated-game conditions so defectors can be punished
        \end{itemize}\item \textbf{Examples}: Nations adopting climate accords; tech industry standardising on USB-C; startups choosing equity over salary; dam-building village cooperatives.
      \end{itemize}
    \end{column}
  \end{columns}
\end{frame}

%%%%%%%%%%%%%%%%%%%%%%%%%%%%%%%%%%%%%%%%%%%%%%%%%%%%%%%%%%%%%%%%%%%%%%%%%%%%%%%%%%
\begin{frame}[fragile]\frametitle{Chicken Game: Ramesh-Suresh 24/7 Standoff}
\begin{columns}
    \begin{column}[T]{0.6\linewidth}
      \textbf{The Scenario:}
      \begin{itemize}
        \item Both decide to stay open 24/7 to outlast each other
        \item Recklessly spending on electricity and overnight staff  
        \item Like two trucks heading toward each other
        \item Winner is whoever has bigger bank account to absorb losses
        \item But if neither blinks, both go bankrupt
        \item Worst outcome: mutual destruction
      \end{itemize}
      
      \vspace{6pt}
      \textbf{Payoff Matrix (Monthly Net Profit/Loss in Rs):}
      \begin{center}
      \begin{tabular}{|c|c|c|}
        \hline
        \textbf{Ramesh $\backslash$ Suresh} & \textbf{Normal Hours} & \textbf{24/7} \\
        \hline
        \textbf{Normal Hours} & (40k, 40k) & (10k, 55k) \\
        \hline
        \textbf{24/7} & (55k, 10k) & (-20k, -20k) \\
        \hline
      \end{tabular}
      \end{center}
    \end{column}
    \begin{column}[T]{0.38\linewidth}
      \textbf{Analysis Steps:}
      \begin{itemize}
        \item \textbf{Check Outcomes:}
        \begin{itemize}
          \item Both normal: Good for both (40K each)
          \item One blinks: Blinker loses badly (10K vs 55K)
          \item Neither blinks: Both lose (Rs 20K)
        \end{itemize}
        \item \textbf{Two Nash Equilibria:} (Normal, 24/7) and (24/7, Normal)
        \item \textbf{Key Difference:} Someone MUST yield
        \item \textbf{Winning Strategy:} Convince opponent you'll never back down
        \item \textbf{Real Danger:} Both commit and crash
      \end{itemize}
    \end{column}
  \end{columns}
\end{frame}


%%%%%%%%%%%%%%%%%%%%%%%%%%%%%%%%%%%%%%%%%%%%%%%%%%%%%%%%%%%%%%%%%%%%%%%%%%%%%%%%%%
\begin{frame}[fragile]\frametitle{Battle of the Sexes: Preference Conflicts}
\begin{columns}
    \begin{column}[T]{0.5\linewidth}
      \begin{itemize}
        \item \textbf{Introduction}: Both players want to coordinate but disagree on \textit{which} coordination point to choose. It is not pure conflict, both want togetherness, but power determines whose preference wins.
        \item \textbf{Situation}: Merger decisions (who becomes CEO?), standard-setting wars (Mac vs PC), coalition government formation, meeting location choices, holiday planning.
      \end{itemize}
    \end{column}
    \begin{column}[T]{0.5\linewidth}
      \begin{itemize}
        \item \textbf{Steps}:
        \begin{itemize}
          \item Confirm both parties genuinely prefer coordination to no deal
          \item Identify the current default, who benefits from the status quo?
          \item Use side-payments or alternating concessions (``I'll watch cricket now, you choose next time'')
          \item Make concessions visible and reciprocal
          \item Formalise agreements to avoid renegotiation
        \end{itemize}\item \textbf{Examples}: Merger CEO selection; choosing a vacation destination; technology standard wars; political coalition ministry allocation.
      \end{itemize}
    \end{column}
  \end{columns}
\end{frame}

%%%%%%%%%%%%%%%%%%%%%%%%%%%%%%%%%%%%%%%%%%%%%%%%%%%%%%%%%%%%%%%%%%%%%%%%%%%%%%%%%%
\begin{frame}[fragile]\frametitle{Stag Hunt: Ramesh-Suresh Quality Upgrade}
\begin{columns}
    \begin{column}[T]{0.59\linewidth}
      \textbf{The Scenario:}
      \begin{itemize}
        \item Both invest in quality beans, train staff, fair prices: street becomes coffee destination (Stag)
        \item Or cut prices, save on quality, steal customers today (Hare)
        \item If one commits to quality while other defects, committed one bears higher costs
        \item Both hesitate, waiting for signal of trust
        \item Challenge: assurance, not conflict
		\item Quality (Stag), Cut Costs (Hare)
      \end{itemize}
      
      \vspace{2pt}
      \textbf{Payoff Matrix (Monthly Profit in Rs):}
      \begin{center}
      \begin{tabular}{|c|c|c|}
        \hline
        \textbf{Ramesh $\backslash$ Suresh} & \textbf{Stag} & \textbf{Hare} \\
        \hline
        \textbf{Stag} & \textbf{(80k, 80k)} & (30k, 55k) \\
        \hline
        \textbf{Hare} & (55k, 30k) & \textbf{(45k, 45k)} \\
        \hline
      \end{tabular}
      \end{center}
    \end{column}
    \begin{column}[T]{0.41\linewidth}
      \textbf{Analysis Steps:}
      \begin{itemize}
        \item \textbf{Two Nash Equilibria:}
        \begin{itemize}
          \item (Quality, Quality): Superior outcome (80K each)
          \item (Cut Costs, Cut Costs): Safe but inferior (45K each)
        \end{itemize}
        \item \textbf{The Problem:} Quality only works if BOTH commit
        \item \textbf{Risk:} If you choose Quality and other defects, you get only 30K
        \item \textbf{Solution:} Build trust, pre-commitment, repeated interaction
        \item \textbf{Key Insight:} About mutual assurance, not defeating opponent
      \end{itemize}
    \end{column}
  \end{columns}
\end{frame}


%%%%%%%%%%%%%%%%%%%%%%%%%%%%%%%%%%%%%%%%%%%%%%%%%%%%%%%%%%%%%%%%%%%%%%%%%%%%%%%%%%
\begin{frame}[fragile]\frametitle{Ultimatum Game: Fairness Over Profit}
\begin{columns}
    \begin{column}[T]{0.5\linewidth}
      \begin{itemize}
        \item \textbf{Introduction}: People will sacrifice their own gain to punish perceived unfairness. This game demonstrates that humans are emotional moralists, not purely rational profit-maximisers, destroying the ``Rational Man'' economic model.
        \item \textbf{Situation}: Any offer/rejection setting: wage negotiations, settlement offers, supplier contracts, brand boycotts. Whenever perceived fairness can override strict economic logic.
      \end{itemize}
    \end{column}
    \begin{column}[T]{0.5\linewidth}
      \begin{itemize}
        \item \textbf{Steps}:
        \begin{itemize}
          \item Frame offers so they appear fair, not just profitable to you
          \item Use anchoring to set fairness expectations early
          \item Give the other party a sense of agency in the process
          \item Avoid visibly large asymmetries in payoffs
          \item Expect rejection of ``logically reasonable'' offers that feel insulting
        \end{itemize}\item \textbf{Examples}: Labour strikes rejecting technically positive wage offers; restaurant tipping even when you will never return; brand boycotts over perceived exploitation.
      \end{itemize}
    \end{column}
  \end{columns}
\end{frame}


%%%%%%%%%%%%%%%%%%%%%%%%%%%%%%%%%%%%%%%%%%%%%%%%%%%%%%%%%%%%%%%%%%%%%%%%%%%%%%%%%%
\begin{frame}[fragile]\frametitle{Ultimatum Game: Ramesh-Suresh Profit Split}
\begin{columns}
    \begin{column}[T]{0.6\linewidth}
      \textbf{The Scenario:}
      \begin{itemize}
        \item Joint catering contract worth Rs 100,000
        \item Ramesh (proposer) must split profit with Suresh (responder)
        \item Ramesh offers a split
        \item Suresh can accept (both get their shares) or reject (both get zero)
        \item Rationally, any positive amount better than zero
        \item But fairness matters, people reject unfair offers
      \end{itemize}
      
      \vspace{2pt}
      \textbf{Possible Outcomes:}
      \begin{center}
      \begin{tabular}{|l|c|c|}
        \hline
        \textbf{Ramesh's Offer} & \textbf{If Accepted} & \textbf{If Rejected} \\
        \hline
        90K:10K split & (90k, 10k) & (0, 0) \\
        \hline
        70K:30K split & (70k, 30k) & (0, 0) \\
        \hline
        50K:50K split & (50k, 50k) & (0, 0) \\
        \hline
      \end{tabular}
      \end{center}
    \end{column}
    \begin{column}[T]{0.4\linewidth}
      \textbf{Analysis:}
      \begin{itemize}
        \item \textbf{Rational Prediction:} Suresh should accept any positive amount (even Rs 1)
        \item \textbf{Reality:} People reject offers below 20-30\%
        \item \textbf{Fairness Norm:} Most offers around 40-50\%
        \item \textbf{Why:} Fairness valued over pure profit
        \item \textbf{Key Insight:} Humans aren't purely profit-maximizing, social norms and emotions matter
        \item \textbf{Strategic Lesson:} Leave enough on table for others
      \end{itemize}
    \end{column}
  \end{columns}
\end{frame}


%%%%%%%%%%%%%%%%%%%%%%%%%%%%%%%%%%%%%%%%%%%%%%%%%%%%%%%%%%%%%%%%%%%%%%%%%%%%%%%%%%
\begin{frame}[fragile]\frametitle{Battle of the Sexes: Ramesh-Suresh Bakery Partnership}
\begin{columns}
    \begin{column}[T]{0.55\linewidth}
      \textbf{The Scenario:}
      \begin{itemize}
        \item Both want ``Breakfast \& Coffee'' combo with local bakery
        \item Ramesh wants bakery to use his premium beans for morning catering
        \item Suresh wants bakery to provide exclusive ``Suresh-branded'' croissants
        \item Both want partnership, both want cooperation
        \item Just can't agree on terms
        \item Struggle: who blinks first and accepts second-best choice?
      \end{itemize}
      
      \vspace{2pt}
      \textbf{Payoff Matrix (Monthly Additions):}
      \begin{center}
      \begin{tabular}{|c|c|c|}
        \hline
        \textbf{Ramesh $\backslash$ Suresh} & \textbf{Ramesh's Terms} & \textbf{Suresh's Terms} \\
        \hline
        \textbf{Ramesh's Terms} & (0, 0) & (30k, 50k) \\
        \hline
        \textbf{Suresh's Terms} & (50k, 30k) & (0, 0) \\
        \hline
      \end{tabular}
      \end{center}
    \end{column}
    \begin{column}[T]{0.45\linewidth}
      \textbf{Analysis Steps:}
      \begin{itemize}
        \item \textbf{Two Nash Equilibria:}
        \begin{itemize}
          \item (Suresh's, Ramesh's): Ramesh gets less but agrees
          \item (Ramesh's, Suresh's): Suresh gets less but agrees
        \end{itemize}
        \item \textbf{Key Feature:} Both prefer coordinating on EITHER outcome over no deal (0, 0)
        \item \textbf{But:} Each prefers different coordination point
        \item \textbf{Power:} With one who cares less, appears more stubborn, or commits first
        \item \textbf{Solutions:} Turn-taking, side payments, pre-commitment
      \end{itemize}
    \end{column}
  \end{columns}
\end{frame}


%%%%%%%%%%%%%%%%%%%%%%%%%%%%%%%%%%%%%%%%%%%%%%%%%%%%%%%%%%%%%%%%%%%%%%%%%%%%%%%%%%
\begin{frame}[fragile]\frametitle{Public Goods Game: The Free-Rider Problem}
\begin{columns}
    \begin{column}[T]{0.5\linewidth}
      \begin{itemize}
        \item \textbf{Introduction}: A public good is non-excludable and non-rivalrous. The free-rider problem emerges when individuals consume more than their fair share while paying less than their fair share, destroying the commons.
        \item \textbf{Situation}: Wikipedia donations, national defence, clean air, office coffee funds, open-source software maintenance, and voluntary neighbourhood upkeep.
      \end{itemize}
    \end{column}
    \begin{column}[T]{0.5\linewidth}
      \begin{itemize}
        \item \textbf{Steps}:
        \begin{itemize}
          \item Identify whether the good is genuinely non-excludable
          \item Introduce social norms or visibility (people contribute more when observed)
          \item Use conditional cooperation mechanisms (``I give if others give'')
          \item Create excludability where possible (paywalls, memberships)
          \item Regulate via taxation or mandated contribution
        \end{itemize}\item \textbf{Examples}: Wikipedia donation drives; national vaccination campaigns; office kitchen cleanliness; funding public broadcasting.
      \end{itemize}
    \end{column}
  \end{columns}
\end{frame}


%%%%%%%%%%%%%%%%%%%%%%%%%%%%%%%%%%%%%%%%%%%%%%%%%%%%%%%%%%%%%%%%%%%%%%%%%%%%%%%%%%
\begin{frame}[fragile]\frametitle{Public Goods: Ramesh-Suresh Street Lighting}
\begin{columns}
    \begin{column}[T]{0.62\linewidth}
      \textbf{The Scenario:}
      \begin{itemize}
        \item Street needs better lighting, costs Rs 20,000
        \item Each can contribute Rs 10,000 or nothing
        \item If installed, BOTH benefit (Rs 15,000 value each)
        \item If one contributes alone, they pay Rs 10K, get Rs 15K benefit = net Rs 5K
        \item If both contribute, each pays Rs 10K, gets Rs 15K = net Rs 5K each
        \item Temptation: let other pay, enjoy benefit for free
      \end{itemize}
      
      \vspace{2pt}
      \textbf{Payoff Matrix (Net Benefit in Rs):}
      \begin{center}
      \begin{tabular}{|c|c|c|}
        \hline
        \textbf{Ramesh $\backslash$ Suresh} & \textbf{Contribute} & \textbf{Free-Ride} \\
        \hline
        \textbf{Contribute} & (5k, 5k) & (-10k, 15k) \\
        \hline
        \textbf{Free-Ride} & (15k, -10k) & \textbf{(0, 0)} \\
        \hline
      \end{tabular}
      \end{center}
    \end{column}
    \begin{column}[T]{0.4\linewidth}
      \textbf{Analysis:}
      \begin{itemize}
        \item \textbf{Dominant Strategy:} Free-ride for both
        \item \textbf{Nash Equilibrium:} (Free-Ride, Free-Ride) = (0, 0)
        \item \textbf{Best Outcome:} Both contribute = (5K, 5K)
        \item \textbf{The Problem:} Individual rationality leads to no public good
        \item \textbf{Solutions:} Mandatory contributions, reputation, repeated interaction, social norms
      \end{itemize}
    \end{column}
  \end{columns}
\end{frame}


%%%%%%%%%%%%%%%%%%%%%%%%%%%%%%%%%%%%%%%%%%%%%%%%%%%%%%%%%%%%%%%%%%%%%%%%%%%%%%%%%%
\begin{frame}[fragile]\frametitle{Matching Pennies: Pure Unpredictability}
\begin{columns}
    \begin{column}[T]{0.5\linewidth}
      \begin{itemize}
        \item \textbf{Introduction}: A purely competitive zero-sum game with no stable pure strategy. Your only winning approach is randomisation, if you have a pattern, the opponent exploits it.
        \item \textbf{Situation}: Penalty kicks; tax audit selection; security screening; military patrol routes; marketing launch timing against competitors, anywhere your adversary benefits from predicting you.
      \end{itemize}
    \end{column}
    \begin{column}[T]{0.5\linewidth}
      \begin{itemize}
        \item \textbf{Steps}:
        \begin{itemize}
          \item Identify that you are in a purely competitive, no-equilibrium-in-pure-strategies game
          \item Calculate the optimal mixed strategy probabilities
          \item Introduce genuine randomisation (dice, random number generators)
          \item Resist the temptation to use ``gut feel'' patterns
          \item Monitor whether the opponent is exploiting any pattern and adjust
        \end{itemize}\item \textbf{Examples}: Penalty kick direction; customs inspection selection; submarine evasion routes; Rock-Paper-Scissors.
      \end{itemize}
    \end{column}
  \end{columns}
\end{frame}


%%%%%%%%%%%%%%%%%%%%%%%%%%%%%%%%%%%%%%%%%%%%%%%%%%%%%%%%%%%%%%%%%%%%%%%%%%%%%%%%%%
\begin{frame}[fragile]\frametitle{Matching Pennies: Ramesh-Suresh Timing Game}
\begin{columns}
    \begin{column}[T]{0.6\linewidth}
      \textbf{The Scenario:}
      \begin{itemize}
        \item Ramesh and Suresh compete for morning rush customers
        \item Each must choose: Open Early (6 AM) or Normal (8 AM)
        \item Ramesh wins if choices match (both early or both normal)
        \item Suresh wins if choices differ (one early, one normal)
        \item Pure conflict, no cooperation possible
        \item Predictability is dangerous
      \end{itemize}
      
      \vspace{6pt}
      \textbf{Payoff Matrix (Customer Advantage):}
      \begin{center}
      \begin{tabular}{|c|c|c|}
        \hline
        \textbf{Ramesh $\backslash$ Suresh} & \textbf{Early} & \textbf{Normal} \\
        \hline
        \textbf{Early} & (1, -1) & (-1, 1) \\
        \hline
        \textbf{Normal} & (-1, 1) & (1, -1) \\
        \hline
      \end{tabular}
      \end{center}
    \end{column}
    \begin{column}[T]{0.38\linewidth}
      \textbf{Analysis:}
      \begin{itemize}
        \item \textbf{No Pure Nash Equilibrium:} Always incentive to deviate
        \item \textbf{Solution:} Mixed strategy, randomize choices
        \item \textbf{Optimal:} Each plays Early 50\%, Normal 50\%
        \item \textbf{Key Insight:} Unpredictability is strategic
        \item \textbf{Examples:} Penalty kicks, rock-paper-scissors, military tactics
        \item \textbf{Lesson:} When opponent can exploit patterns, randomize
      \end{itemize}
    \end{column}
  \end{columns}
\end{frame}


%%%%%%%%%%%%%%%%%%%%%%%%%%%%%%%%%%%%%%%%%%%%%%%%%%%%%%%%%%%%%%%%%%%%%%%%%%%%%%%%%%
\begin{frame}[fragile]\frametitle{Cooperative Game Theory: Sharing the Gains}
\begin{columns}
    \begin{column}[T]{0.5\linewidth}
      \begin{itemize}
        \item \textbf{Introduction}: Cooperative game theory asks: when players can form coalitions and make binding agreements, how should gains be fairly and stably divided? The Shapley Value and the Core are the key solution concepts.
        \item \textbf{Situation}: Startup profit-sharing, political coalition formation, airline alliances, pharmaceutical joint ventures, supply-chain cost allocation, international carbon agreements.
      \end{itemize}
    \end{column}
    \begin{column}[T]{0.5\linewidth}
      \begin{itemize}
        \item \textbf{Steps}:
        \begin{itemize}
          \item Identify all possible coalitions and their total values
          \item Calculate each player's marginal contribution (Shapley Value)
          \item Check if the allocation is in the Core, no sub-coalition wants to break away
          \item Negotiate side-payments if the allocation is outside the core
          \item Build binding agreements and enforcement mechanisms
        \end{itemize}\item \textbf{Examples}: Law firm profit splits; coalition government ministry allocations; airline revenue-sharing alliances; joint-venture R\&D cost-sharing.
      \end{itemize}
    \end{column}
  \end{columns}
\end{frame}


%%%%%%%%%%%%%%%%%%%%%%%%%%%%%%%%%%%%%%%%%%%%%%%%%%%%%%%%%%%%%%%%%%%%%%%%%%%%%%%%%%
\begin{frame}[fragile]\frametitle{Cooperative Game Theory: Ramesh-Suresh Joint Venture}
\begin{columns}
    \begin{column}[T]{0.65\linewidth}
      \textbf{The Scenario:}
      \begin{itemize}
        \item Both can bid jointly for corporate catering contract
        \item Alone: Ramesh can earn Rs 100K, Suresh Rs 80K
        \item Together: Combined capacity = Rs 250K total
        \item Question: How to split fairly?
        \item Shapley Value: Based on marginal contribution
        \item The Core: Allocations where no one wants to leave
      \end{itemize}
      
      \vspace{2pt}
      \textbf{Value Analysis:}
      \begin{center}
      \small
      \begin{tabular}{|l|c|}
        \hline
        \textbf{Coalition} & \textbf{Total Value} \\
        \hline
        Ramesh alone & Rs 100K \\
        Suresh alone & Rs 80K \\
        Both together & Rs 250K \\
        \hline
        \textbf{Surplus from cooperation} & \textbf{Rs 70K} \\
        \hline
      \end{tabular}
      \end{center}
    \end{column}
    \begin{column}[T]{0.35\linewidth}
      \textbf{Fair Split Approaches:}
      \begin{itemize}
        \item \textbf{Shapley Value:}
        \begin{itemize}
          \item Ramesh: Rs 135K
          \item Suresh: Rs 115K
        \end{itemize}
        \item \textbf{Must be in Core:}
        \begin{itemize}
          \item Ramesh gets $\geq$ Rs 100K
          \item Suresh gets $\geq$ Rs 80K
          \item Sum = Rs 250K
        \end{itemize}
        \item \textbf{Key Insight:} Fair split considers outside options and marginal contributions
        \item \textbf{Applications:} Startups, coalitions, partnerships
      \end{itemize}
    \end{column}
  \end{columns}
\end{frame}


% %%%%%%%%%%%%%%%%%%%%%%%%%%%%%%%%%%%%%%%%%%%%%%%%%%%%%%%%%%%%%%%%%%%%%%%%%%%%%%%%%%
% \begin{frame}[fragile]\frametitle{}
% \begin{center}
% {\Large Decision-Making Approaches}
% \end{center}
% \end{frame}

%%%%%%%%%%%%%%%%%%%%%%%%%%%%%%%%%%%%%%%%%%%%%%%%%%%%%%%%%%%%%%%%%%%%%%%%%%%%%%%%%%
\begin{frame}[fragile]\frametitle{Backward Induction: Think Backwards, Act Forward}
\begin{columns}
    \begin{column}[T]{0.5\linewidth}
      \begin{itemize}
        \item \textbf{Introduction}: To solve a sequential game, examine the final stage, find the optimal move there, then roll back logic to the start. It prevents moves today that create traps tomorrow.
        \item \textbf{Situation}: Chess endgames, business exit strategy design, career planning, negotiations with defined endpoints, syllabus design, retirement planning.
      \end{itemize}
    \end{column}
    \begin{column}[T]{0.5\linewidth}
      \begin{itemize}
        \item \textbf{Steps}:
        \begin{itemize}
          \item Define the final desired outcome clearly
          \item Work backwards: what must be true one step before the end?
          \item Continue rolling back, mapping each decision node
          \item Identify the first move that leads to the desired end
          \item Check for credibility, are all intermediate steps actually incentive-compatible?
        \end{itemize}\item \textbf{Examples}: Chess endgame analysis; Indian parents planning children's education backwards from career goal; startup exit strategy mapped from IPO back to MVP.
      \end{itemize}
    \end{column}
  \end{columns}
\end{frame}


%%%%%%%%%%%%%%%%%%%%%%%%%%%%%%%%%%%%%%%%%%%%%%%%%%%%%%%%%%%%%%%%%%%%%%%%%%%%%%%%%%
\begin{frame}[fragile]\frametitle{Backward Induction: Ramesh Plans His Exit}
\begin{columns}
    \begin{column}[T]{0.6\linewidth}
      \textbf{The Scenario:}
      \begin{itemize}
        \item Ramesh plans to sell his shop in 3 years
        \item Starts by imagining: What will buyer want?
        \item Answer: Stable brand, loyal customers, healthy margins
        \item Works backward: Price war today would create discount-chasing customers
        \item Brand would be worthless at exit
        \item So today's move: Raise quality, redesign space, position as premium
      \end{itemize}
      
      \vspace{2pt}
      \textbf{Decision Tree (Simplified):}
      \begin{center}
      \small
      \begin{tabular}{|l|l|c|}
        \hline
        \textbf{Today} & \textbf{In 3 Years} & \textbf{Sale Value} \\
        \hline
        Price War & Discount  & Low (Rs 10L) \\
        \hline
        Premium Quality & Loyal high-value  & High (Rs 50L) \\
        \hline
      \end{tabular}
      \end{center}
    \end{column}
    \begin{column}[T]{0.38\linewidth}
      \textbf{Steps:}
      \begin{itemize}
        \item \textbf{Define End Goal:} Sell for maximum value
        \item \textbf{Work Backward:} What creates value at exit? Loyal premium customers
        \item \textbf{One Step Before:} What builds loyalty? Consistent quality
        \item \textbf{Earlier Still:} What enables quality? Investment today
        \item \textbf{First Move:} Raise quality now, even if some customers leave
        \item \textbf{Key Insight:} Today's "loss" is tomorrow's setup
      \end{itemize}
    \end{column}
  \end{columns}
\end{frame}


%%%%%%%%%%%%%%%%%%%%%%%%%%%%%%%%%%%%%%%%%%%%%%%%%%%%%%%%%%%%%%%%%%%%%%%%%%%%%%%%%%
\begin{frame}[fragile]\frametitle{Forward Induction: Reading Hidden Intentions}
\begin{columns}
    \begin{column}[T]{0.5\linewidth}
      \begin{itemize}
        \item \textbf{Introduction}: Past actions reveal future intent. If a player makes a move that would be irrational \textit{unless} they held a specific strong position, you can infer that strong position.
        \item \textbf{Situation}: Interpreting competitor signals in markets, reading investor confidence from founder behaviour, recruiting and interviewing, diplomacy, strategic partnership evaluation.
      \end{itemize}
    \end{column}
    \begin{column}[T]{0.5\linewidth}
      \begin{itemize}
        \item \textbf{Steps}:
        \begin{itemize}
          \item Observe a costly or unusual action by the other party
          \item Ask: what beliefs would make this action rational?
          \item Update your assessment of their type or intention
          \item Adjust your strategy accordingly
          \item Use the same logic to send credible signals about your own type
        \end{itemize}\item \textbf{Examples}: A founder taking a Rs. 1 salary signals supreme confidence to investors; large Super Bowl ad spend signals financial strength; university degrees signal perseverance.
      \end{itemize}
    \end{column}
  \end{columns}
\end{frame}


%%%%%%%%%%%%%%%%%%%%%%%%%%%%%%%%%%%%%%%%%%%%%%%%%%%%%%%%%%%%%%%%%%%%%%%%%%%%%%%%%%
\begin{frame}[fragile]\frametitle{Forward Induction: Ramesh's Price Freeze Signal}
\begin{columns}
    \begin{column}[T]{0.6\linewidth}
      \textbf{The Scenario:}
      \begin{itemize}
        \item Ramesh announces: Prices unchanged for 3 years, despite rising costs
        \item Looks reckless to Suresh
        \item But Ramesh is burning flexibility to signal confidence
        \item If Ramesh were unsure about footfall, this would box him into losses
        \item The fact he accepts risk suggests hidden strength
        \item Maybe superior location, loyal base, or cost advantage
      \end{itemize}
      
      \vspace{6pt}
      \textbf{Suresh's Inference:}
      \begin{center}
      \small
      \begin{tabular}{|l|l|}
        \hline
        \textbf{Ramesh's Move} & \textbf{What It Signals} \\
        \hline
        Price freeze & Confident in volume \\
        Accepts margin pressure & Has cost advantage \\
        Public commitment & Won't back down \\
        \hline
      \end{tabular}
      \end{center}
    \end{column}
    \begin{column}[T]{0.38\linewidth}
      \textbf{Analysis Steps:}
      \begin{itemize}
        \item \textbf{Observe Costly Action:} Price freeze with rising costs
        \item \textbf{Ask:} What beliefs make this rational?
        \item \textbf{Answer:} Confidence in strong customer base
        \item \textbf{Infer Type:} Ramesh likely has advantage
        \item \textbf{Update Strategy:} Don't compete on price, differentiate elsewhere
        \item \textbf{Key Insight:} Actions speak louder than words, costly signals are credible
      \end{itemize}
    \end{column}
  \end{columns}
\end{frame}


%%%%%%%%%%%%%%%%%%%%%%%%%%%%%%%%%%%%%%%%%%%%%%%%%%%%%%%%%%%%%%%%%%%%%%%%%%%%%%%%%%
\begin{frame}[fragile]\frametitle{Sequential Games: Turn-Based Strategy}
\begin{columns}
    \begin{column}[T]{0.5\linewidth}
      \begin{itemize}
        \item \textbf{Introduction}: In sequential games, players move one after another and can observe prior moves. Time and order of moves determine power, first movers set the tone, last movers often control the outcome.
        \item \textbf{Situation}: Real estate negotiation, chess, lobbying, wage bargaining, courtroom strategy, Stackelberg competition in duopoly markets.
      \end{itemize}
    \end{column}
    \begin{column}[T]{0.5\linewidth}
      \begin{itemize}
        \item \textbf{Steps}:
        \begin{itemize}
          \item Map the game tree with all decision nodes
          \item Apply backward induction
          \item Identify who benefits from moving first vs.\ last
          \item Consider pre-commitment to shape subsequent play
          \item Use observable actions to constrain the opponent's future choices
        \end{itemize}\item \textbf{Examples}: House purchase negotiation (offer $\rightarrow$ counter $\rightarrow$ walk away $\rightarrow$ call back); Stackelberg leader in output competition; chess opening theory.
      \end{itemize}
    \end{column}
  \end{columns}
\end{frame}


%%%%%%%%%%%%%%%%%%%%%%%%%%%%%%%%%%%%%%%%%%%%%%%%%%%%%%%%%%%%%%%%%%%%%%%%%%%%%%%%%%
\begin{frame}[fragile]\frametitle{Sequential Games: Ramesh-Suresh Second Floor}
\begin{columns}
    \begin{column}[T]{0.6\linewidth}
      \textbf{The Scenario:}
      \begin{itemize}
        \item Move 1: Ramesh adds second floor for remote workers
        \item Move 2: Suresh observes, then introduces quiet zone + faster WiFi
        \item Move 3: Ramesh notices noise complaints, adds premium pricing
        \item Move 4: Suresh offers membership passes with unlimited refills
        \item Move 5: Ramesh pivots to evening networking events
        \item Move 6: Suresh extends hours, markets as all-day workspace
      \end{itemize}
      
      \vspace{4pt}
      \textbf{Key Feature: Observable Moves}
      \begin{center}
      \small Each player sees previous moves before deciding
      \end{center}
    \end{column}
    \begin{column}[T]{0.38\linewidth}
      \textbf{Strategic Advantages:}
      \begin{itemize}
        \item \textbf{First Mover:} 
        \begin{itemize}
          \item Frames the game
          \item Signals commitment
          \item Sets agenda
        \end{itemize}
        \item \textbf{Second Mover:}
        \begin{itemize}
          \item Gains information
          \item Can tailor response
          \item Often controls outcome
        \end{itemize}
        \item \textbf{Key Tools:} Backward induction, commitment, timing
        \item \textbf{Real Examples:} Negotiations, lobbying, chess
      \end{itemize}
    \end{column}
  \end{columns}
\end{frame}


%%%%%%%%%%%%%%%%%%%%%%%%%%%%%%%%%%%%%%%%%%%%%%%%%%%%%%%%%%%%%%%%%%%%%%%%%%%%%%%%%%
\begin{frame}[fragile]\frametitle{Simultaneous Games: Deciding Under Uncertainty}
\begin{columns}
    \begin{column}[T]{0.5\linewidth}
      \begin{itemize}
        \item \textbf{Introduction}: Players act at the same time without observing each other's choices. Strategy must be based on probability and beliefs about opponent rationality, not observed moves.
        \item \textbf{Situation}: Sealed tender auctions, stock market investing, general elections, simultaneous product launches, penalty kicks where the decision is truly simultaneous.
      \end{itemize}
    \end{column}
    \begin{column}[T]{0.5\linewidth}
      \begin{itemize}
        \item \textbf{Steps}:
        \begin{itemize}
          \item Enumerate all strategies and build a payoff matrix
          \item Search for dominant strategies
          \item If none exist, find Nash Equilibrium in pure strategies
          \item If none in pure strategies, compute mixed strategy Nash Equilibrium
          \item Use prior information about opponent types to refine predictions
        \end{itemize}\item \textbf{Examples}: Sealed government contract bids; buying stocks before earnings announcements; Rock-Paper-Scissors.
      \end{itemize}
    \end{column}
  \end{columns}
\end{frame}


%%%%%%%%%%%%%%%%%%%%%%%%%%%%%%%%%%%%%%%%%%%%%%%%%%%%%%%%%%%%%%%%%%%%%%%%%%%%%%%%%%
\begin{frame}[fragile]\frametitle{Simultaneous Games: Ramesh-Suresh Menu Launch}
\begin{columns}
    \begin{column}[T]{0.6\linewidth}
      \textbf{The Scenario:}
      \begin{itemize}
        \item Both planning to launch new menu items this month
        \item Must decide without knowing other's choice
        \item Options: Italian (pasta/pizza) or Asian (noodles/rice)
        \item If different: Each captures niche market
        \item If same: Split market, higher competition
        \item Cannot observe each other's prep until launch day
      \end{itemize}
      
      \vspace{2pt}
      \textbf{Payoff Matrix (Additional Revenue/m):}
      \begin{center}
      \begin{tabular}{|c|c|c|}
        \hline
        \textbf{Ramesh $\backslash$ Suresh} & \textbf{Italian} & \textbf{Asian} \\
        \hline
        \textbf{Italian} & (20k, 20k) & (40k, 40k) \\
        \hline
        \textbf{Asian} & (40k, 40k) & (20k, 20k) \\
        \hline
      \end{tabular}
      \end{center}
    \end{column}
    \begin{column}[T]{0.38\linewidth}
      \textbf{Analysis:}
      \begin{itemize}
        \item \textbf{No Dominant Strategy:} Best choice depends on opponent
        \item \textbf{Two Nash Equilibria:} (Italian, Asian) and (Asian, Italian)
        \item \textbf{Challenge:} How to coordinate without communication?
        \item \textbf{Solutions:}
        \begin{itemize}
          \item Past patterns
          \item Focal points
          \item Market research
          \item Mixed strategies
        \end{itemize}
        \item \textbf{Key Difference:} Must decide blindly, no info about opponent's move
      \end{itemize}
    \end{column}
  \end{columns}
\end{frame}


%%%%%%%%%%%%%%%%%%%%%%%%%%%%%%%%%%%%%%%%%%%%%%%%%%%%%%%%%%%%%%%%%%%%%%%%%%%%%%%%%%
\begin{frame}[fragile]\frametitle{Subgame Perfect Equilibrium: Credible Plans}
\begin{columns}
    \begin{column}[T]{0.5\linewidth}
      \begin{itemize}
        \item \textbf{Introduction}: Nash equilibrium can include non-credible threats that are never tested. Subgame perfection demands that your strategy must be optimal at \textit{every} decision point, even in branches of the game tree you hope never to reach.
        \item \textbf{Situation}: Labor negotiations, military deterrence, corporate takeover defence, parent-child discipline, contract enforcement, anywhere threats or promises must be believed.
      \end{itemize}
    \end{column}
    \begin{column}[T]{0.5\linewidth}
      \begin{itemize}
        \item \textbf{Steps}:
        \begin{itemize}
          \item Draw the full game tree
          \item Apply backward induction at every node
          \item Remove strategies that are irrational at any sub-game
          \item Ask: if this branch were reached, would I actually follow through?
          \item If not, redesign your commitment to make it credible
        \end{itemize}\item \textbf{Examples}: Parent who always carries through on stated consequences; proportional military response doctrines; poison-pill takeover defences that are actually executable.
      \end{itemize}
    \end{column}
  \end{columns}
\end{frame}

%%%%%%%%%%%%%%%%%%%%%%%%%%%%%%%%%%%%%%%%%%%%%%%%%%%%%%%%%%%%%%%%%%%%%%%%%%%%%%%%%%
\begin{frame}[fragile]\frametitle{Minimax Strategy: Minimising Maximum Loss}
\begin{columns}
    \begin{column}[T]{0.5\linewidth}
      \begin{itemize}
        \item \textbf{Introduction}: A conservative strategy that looks at the worst-case scenario of every decision and selects the one with the least-bad maximum loss. It is playing not to lose rather than playing to win.
        \item \textbf{Situation}: Risk management, insurance, cybersecurity, crisis planning, healthcare protocols, wherever catastrophic downside must be avoided even at cost of upside.
      \end{itemize}
    \end{column}
    \begin{column}[T]{0.5\linewidth}
      \begin{itemize}
        \item \textbf{Steps}:
        \begin{itemize}
          \item List all your strategies and all opponent strategies
          \item Find the worst payoff for each of your strategies (the ``minimum'')
          \item Select the strategy with the highest minimum payoff
          \item Compare minimax cost against expected value analysis
          \item Adjust based on risk tolerance and reversibility of outcomes
        \end{itemize}\item \textbf{Examples}: Travel insurance; diversified investment portfolios; pre-nuptial agreements; cybersecurity breach response planning.
      \end{itemize}
    \end{column}
  \end{columns}
\end{frame}

%%%%%%%%%%%%%%%%%%%%%%%%%%%%%%%%%%%%%%%%%%%%%%%%%%%%%%%%%%%%%%%%%%%%%%%%%%%%%%%%%%
\begin{frame}[fragile]\frametitle{Bayesian Games: Decisions with Incomplete Information}
\begin{columns}
    \begin{column}[T]{0.5\linewidth}
      \begin{itemize}
        \item \textbf{Introduction}: Most real-world games involve uncertainty about opponents' types, payoffs, or capabilities. Players form beliefs, observe actions, update via Bayes' rule, and choose strategies maximising expected payoffs given current beliefs.
        \item \textbf{Situation}: Used car markets, job interviews, insurance underwriting, military intelligence assessment, product demand forecasting, dating and relationship formation.
      \end{itemize}
    \end{column}
    \begin{column}[T]{0.5\linewidth}
      \begin{itemize}
        \item \textbf{Steps}:
        \begin{itemize}
          \item Assign prior probability distributions over opponent types
          \item Identify what signals or actions reveal information
          \item Update beliefs using Bayes' Rule as evidence arrives
          \item Compute expected payoffs across opponent types
          \item Choose strategy maximising expected payoff; recognise your actions also signal your type
        \end{itemize}\item \textbf{Examples}: Employer updating beliefs about candidate quality through interview; used-car buyer reading seller signals; poker player updating hand estimate from betting patterns.
      \end{itemize}
    \end{column}
  \end{columns}
\end{frame}

% %%%%%%%%%%%%%%%%%%%%%%%%%%%%%%%%%%%%%%%%%%%%%%%%%%%%%%%%%%%%%%%%%%%%%%%%%%%%%%%%%%
% \begin{frame}[fragile]\frametitle{}
% \begin{center}
% {\Large Credibility and Commitment}
% \end{center}
% \end{frame}

%%%%%%%%%%%%%%%%%%%%%%%%%%%%%%%%%%%%%%%%%%%%%%%%%%%%%%%%%%%%%%%%%%%%%%%%%%%%%%%%%%
\begin{frame}[fragile]\frametitle{Strategic Commitment: Limiting Your Own Options}
\begin{columns}
    \begin{column}[T]{0.5\linewidth}
      \begin{itemize}
        \item \textbf{Introduction}: Paradoxically, you can increase bargaining power by \textit{reducing} your own options. If you credibly cannot back down, the other party must adapt to your position.
        \item \textbf{Situation}: Wage negotiations, competitive market entry threats, diplomatic standoffs, self-control problems, anywhere restricting future flexibility strengthens your current position.
      \end{itemize}
    \end{column}
    \begin{column}[T]{0.5\linewidth}
      \begin{itemize}
        \item \textbf{Steps}:
        \begin{itemize}
          \item Identify the decision you want the other party to take
          \item Find a way to make your own retreat genuinely costly or impossible
          \item Make the commitment visible and verifiable
          \item Ensure the commitment is specific, not vague
          \item Time the commitment before negotiations begin
        \end{itemize}\item \textbf{Examples}: Cortés burning his ships; automatic pension deductions; publicly announced deadlines; contracts with heavy penalty clauses; hunger strikes.
      \end{itemize}
    \end{column}
  \end{columns}
\end{frame}


%%%%%%%%%%%%%%%%%%%%%%%%%%%%%%%%%%%%%%%%%%%%%%%%%%%%%%%%%%%%%%%%%%%%%%%%%%%%%%%%%%
\begin{frame}[fragile]\frametitle{Strategic Commitment: Suresh's Organic Contract}
\begin{columns}
    \begin{column}[T]{0.6\linewidth}
      \textbf{The Scenario:}
      \begin{itemize}
        \item Suresh tired of endless price wars draining margins
        \item Signs 5-year non-breakable contract with organic bean supplier
        \item Locks into higher input costs
        \item Frames contract copy near counter where all can see
        \item Message to Ramesh: Pricing at Rs 50 now impossible
        \item Removed own escape route to gain leverage
      \end{itemize}
      
      \vspace{6pt}
      \textbf{Before vs After Commitment:}
      \begin{center}
      \small
      \begin{tabular}{|l|l|l|}
        \hline
        \textbf{Scenario} & \textbf{Before} & \textbf{After} \\
        \hline
        Can lower prices? & Yes & No (contract) \\
        Ramesh's threat & Credible & Not credible \\
        Price war & Continues & Ends \\
        \hline
      \end{tabular}
      \end{center}
    \end{column}
    \begin{column}[T]{0.38\linewidth}
      \textbf{Analysis:}
      \begin{itemize}
        \item \textbf{Paradox:} Strength from constraint
        \item \textbf{By giving up flexibility:} Gains influence
        \item \textbf{Key Requirements:}
        \begin{itemize}
          \item Visible (framed contract)
          \item Credible (legally binding)
          \item Costly to reverse
        \end{itemize}
        \item \textbf{Outcome:} Ramesh stops price war, focuses on differentiation
        \item \textbf{Lesson:} Your lack of options becomes their problem
      \end{itemize}
    \end{column}
  \end{columns}
\end{frame}


%%%%%%%%%%%%%%%%%%%%%%%%%%%%%%%%%%%%%%%%%%%%%%%%%%%%%%%%%%%%%%%%%%%%%%%%%%%%%%%%%%
\begin{frame}[fragile]\frametitle{Pre-Commitment Devices: Binding Your Future Self}
\begin{columns}
    \begin{column}[T]{0.5\linewidth}
      \begin{itemize}
        \item \textbf{Introduction}: We are often two people, the planner and the impulsive doer. Pre-commitment devices are strategies the planner uses to constrain the doer before temptation strikes.
        \item \textbf{Situation}: Dieting, saving, avoiding addictions, maintaining study habits, managing social media overuse, any area where your future self's impulsive decisions conflict with your present-self's goals.
      \end{itemize}
    \end{column}
    \begin{column}[T]{0.5\linewidth}
      \begin{itemize}
        \item \textbf{Steps}:
        \begin{itemize}
          \item Identify the specific moment of temptation
          \item Introduce a barrier before that moment (not after)
          \item Make the barrier automatic, not dependent on future willpower
          \item Increase the cost of overriding the commitment
          \item Use social accountability to amplify the pre-commitment
        \end{itemize}\item \textbf{Examples}: Odysseus tying himself to the mast; deleting social media apps; freezing credit cards in ice; gym buddies; automatic savings transfers.
      \end{itemize}
    \end{column}
  \end{columns}
\end{frame}

%%%%%%%%%%%%%%%%%%%%%%%%%%%%%%%%%%%%%%%%%%%%%%%%%%%%%%%%%%%%%%%%%%%%%%%%%%%%%%%%%%
\begin{frame}[fragile]\frametitle{Credible Threats: Making Promises Believable}
\begin{columns}
    \begin{column}[T]{0.5\linewidth}
      \begin{itemize}
        \item \textbf{Introduction}: A threat is only useful when it is credible. If carrying out the threat hurts you more than the opponent, it will be called as a bluff. The incentive structure must make follow-through rational.
        \item \textbf{Situation}: Parental discipline, deterrence strategy, price-match guarantees, union strike threats, contractual penalty clauses, wherever compliance requires the other party to believe punishment is real.
      \end{itemize}
    \end{column}
    \begin{column}[T]{0.5\linewidth}
      \begin{itemize}
        \item \textbf{Steps}:
        \begin{itemize}
          \item Ask honestly: would I actually carry this out?
          \item Adjust the threat so it costs you less than compliance costs the opponent
          \item Make the threat automatic (remove discretion)
          \item Build a track record of following through
          \item Communicate the threat early and unambiguously
        \end{itemize}\item \textbf{Examples}: Nuclear deterrence structures; price-match guarantees (credible threat to rivals); central bank inflation targets; unconditional resignation letters.
      \end{itemize}
    \end{column}
  \end{columns}
\end{frame}

%%%%%%%%%%%%%%%%%%%%%%%%%%%%%%%%%%%%%%%%%%%%%%%%%%%%%%%%%%%%%%%%%%%%%%%%%%%%%%%%%%
\begin{frame}[fragile]\frametitle{Brinkmanship: Playing with Fire}
\begin{columns}
    \begin{column}[T]{0.5\linewidth}
      \begin{itemize}
        \item \textbf{Introduction}: Deliberately pushing a situation to the edge of mutual disaster, or creating the appearance of losing control, to frighten the opponent into conceding. The shared risk becomes the leverage.
        \item \textbf{Situation}: Debt-ceiling negotiations, strikes timed to peak season, geopolitical confrontations, diplomatic ultimatums, late-stage acquisition negotiations.
      \end{itemize}
    \end{column}
    \begin{column}[T]{0.5\linewidth}
      \begin{itemize}
        \item \textbf{Steps}:
        \begin{itemize}
          \item Assess the cost of mutual disaster, it must be roughly equal or worse for the opponent
          \item Create visible risk of escalation (not just the threat)
          \item Maintain plausible deniability, ``I may not be able to stop this.''
          \item Keep an offramp for the opponent to retreat without total humiliation
          \item Know precisely when to de-escalate
        \end{itemize}\item \textbf{Examples}: Indo-Pak border tensions; union strikes during harvest season; debt-ceiling confrontations; Cold War nuclear postures.
      \end{itemize}
    \end{column}
  \end{columns}
\end{frame}

%%%%%%%%%%%%%%%%%%%%%%%%%%%%%%%%%%%%%%%%%%%%%%%%%%%%%%%%%%%%%%%%%%%%%%%%%%%%%%%%%%
\begin{frame}[fragile]\frametitle{Burning Bridges: Eliminating Retreat}
\begin{columns}
    \begin{column}[T]{0.5\linewidth}
      \begin{itemize}
        \item \textbf{Introduction}: By deliberately eliminating your own escape routes, you signal unshakeable commitment. When retreat is impossible, threats become credible, bluffs become real, and opponents must take your position seriously.
        \item \textbf{Situation}: Market entry with irreversible investment, public goal announcements, constitutional lock-ins, military positions, start-up pivots where the old product is shut down.
      \end{itemize}
    \end{column}
    \begin{column}[T]{0.5\linewidth}
      \begin{itemize}
        \item \textbf{Steps}:
        \begin{itemize}
          \item Identify what escape routes the opponent expects you to take
          \item Eliminate those routes publicly and irreversibly
          \item Signal the elimination to your opponent clearly
          \item Ensure the commitment is specific enough to be believed
          \item Carefully weigh: are you certain enough to make retreat impossible?
        \end{itemize}\item \textbf{Examples}: Cortés burning ships in Mexico; no-refund advance-payment policies; publicly announced deadlines; constitutional supermajority requirements; military forces holding indefensible positions.
      \end{itemize}
    \end{column}
  \end{columns}
\end{frame}

%%%%%%%%%%%%%%%%%%%%%%%%%%%%%%%%%%%%%%%%%%%%%%%%%%%%%%%%%%%%%%%%%%%%%%%%%%%%%%%%%%
\begin{frame}[fragile]\frametitle{Reputation Effects: Your History Precedes You}
\begin{columns}
    \begin{column}[T]{0.5\linewidth}
      \begin{itemize}
        \item \textbf{Introduction}: In repeated interactions, your past actions determine your future payoffs. A reputation for being ``tough'' or ``honest'' changes how others play against you, it is a capital asset that earns compound returns.
        \item \textbf{Situation}: Brand management, credit markets, diplomatic relations, business partnerships, personal professional networks, any context with repeated interaction and observable history.
      \end{itemize}
    \end{column}
    \begin{column}[T]{0.5\linewidth}
      \begin{itemize}
        \item \textbf{Steps}:
        \begin{itemize}
          \item Identify what reputation you want to build
          \item Make short-term sacrifices to establish credibility
          \item Ensure your actions are observable
          \item Maintain consistency, reputation is destroyed by exceptions
          \item Invest in platforms where reputation is recorded (ratings, testimonials, public records)
        \end{itemize}\item \textbf{Examples}: Tata Group's century-long reputation; credit scores; eBay seller ratings; central bank inflation credibility; Mafia omertà.
      \end{itemize}
    \end{column}
  \end{columns}
\end{frame}

%%%%%%%%%%%%%%%%%%%%%%%%%%%%%%%%%%%%%%%%%%%%%%%%%%%%%%%%%%%%%%%%%%%%%%%%%%%%%%%%%%
\begin{frame}[fragile]\frametitle{Escalation of Commitment: The Sunk Cost Trap}
\begin{columns}
    \begin{column}[T]{0.5\linewidth}
      \begin{itemize}
        \item \textbf{Introduction}: The irrational decision to continue a failing course of action because of past investments. Rational strategy demands looking only forward, past costs are gone and irrelevant to future decisions.
        \item \textbf{Situation}: Ongoing projects with poor ROI, failing relationships, bad stocks held too long, unfinished degrees, prolonged wars, any situation where sunk cost drives continued investment.
      \end{itemize}
    \end{column}
    \begin{column}[T]{0.5\linewidth}
      \begin{itemize}
        \item \textbf{Steps}:
        \begin{itemize}
          \item Separate past costs from future costs and benefits
          \item Ask: if I were starting today, would I make this investment?
          \item Establish clear kill-criteria in advance
          \item Assign an independent reviewer not emotionally attached to the project
          \item Reframe exit as strategic intelligence, not failure
        \end{itemize}\item \textbf{Examples}: The Vietnam War; waiting for a bus already 30 minutes late; gambling to ``win it back''; keeping a failing business unit alive.
      \end{itemize}
    \end{column}
  \end{columns}
\end{frame}

% %%%%%%%%%%%%%%%%%%%%%%%%%%%%%%%%%%%%%%%%%%%%%%%%%%%%%%%%%%%%%%%%%%%%%%%%%%%%%%%%%%
% \begin{frame}[fragile]\frametitle{}
% \begin{center}
% {\Large Strategic Interaction in Relationships}
% \end{center}
% \end{frame}

%%%%%%%%%%%%%%%%%%%%%%%%%%%%%%%%%%%%%%%%%%%%%%%%%%%%%%%%%%%%%%%%%%%%%%%%%%%%%%%%%%
\begin{frame}[fragile]\frametitle{Tit-for-Tat: Reciprocity Works}
\begin{columns}
    \begin{column}[T]{0.5\linewidth}
      \begin{itemize}
        \item \textbf{Introduction}: The winning strategy in repeated Prisoner's Dilemma tournaments: cooperate first, then mirror whatever the opponent did last round. Simple, forgiving, but retaliatory, the algorithm of human morality.
        \item \textbf{Situation}: Trade negotiations, business partnerships, gifting cultures, arms control agreements, neighbourhood relations, any long-term repeated interaction where defection must be deterred.
      \end{itemize}
    \end{column}
    \begin{column}[T]{0.5\linewidth}
      \begin{itemize}
        \item \textbf{Steps}:
        \begin{itemize}
          \item Open by cooperating unconditionally
          \item Immediately mirror any defection with one round of defection
          \item Return to cooperation the moment the opponent cooperates
          \item Never defect first
          \item Communicate the strategy so the opponent understands the pattern
        \end{itemize}\item \textbf{Examples}: Trade tariff reciprocity; reciprocal dinner invitations; ceasefire agreements conditioned on compliance; eye-for-an-eye justice systems.
      \end{itemize}
    \end{column}
  \end{columns}
\end{frame}


%%%%%%%%%%%%%%%%%%%%%%%%%%%%%%%%%%%%%%%%%%%%%%%%%%%%%%%%%%%%%%%%%%%%%%%%%%%%%%%%%%
\begin{frame}[fragile]\frametitle{Tit-for-Tat: Ramesh-Suresh Reciprocity}
\begin{columns}
    \begin{column}[T]{0.63\linewidth}
      \textbf{The Scenario:}
      \begin{itemize}
        \item Both trying to undercut with lower prices, aggressive promotions
        \item Pattern: Aggression met with aggression, both lose margins
        \item Ramesh tries different: Offers free pastry with morning coffee
        \item Suresh notices, responds with small gesture
        \item Respect boundaries, retaliate only when needed, forgive small moves
        \item Coffee war → cooperation w tit-for-tat
      \end{itemize}
      
      % \vspace{2pt}
      \textbf{Tit-for-Tat Strategy:}
      \begin{center}
      \small
      \begin{tabular}{|l|l|}
        \hline
        \textbf{Step} & \textbf{Action} \\
        \hline
        1. First move & Cooperate (be nice first) \\
        2. If opponent cooperates & Continue cooperating \\
        3. If opponent defects & Retaliate once (proportional) \\
        4. If opponent returns & Forgive, resume cooperation \\
        \hline
      \end{tabular}
      \end{center}
    \end{column}
    \begin{column}[T]{0.37\linewidth}
      \textbf{Why It Works:}
      \begin{itemize}
        \item \textbf{Nice:} Never defects first
        \item \textbf{Retaliatory:} Punishes defection immediately
        \item \textbf{Forgiving:} Returns to cooperation
        \item \textbf{Clear:} Opponent understands pattern
        \item \textbf{Examples:}
        \begin{itemize}
          \item Trade tariffs
          \item Social invitations
          \item Ceasefire agreements
        \end{itemize}
        \item \textbf{Key Insight:} Golden Rule with teeth, cooperation can survive among self-interested agents
      \end{itemize}
    \end{column}
  \end{columns}
\end{frame}


%%%%%%%%%%%%%%%%%%%%%%%%%%%%%%%%%%%%%%%%%%%%%%%%%%%%%%%%%%%%%%%%%%%%%%%%%%%%%%%%%%
\begin{frame}[fragile]\frametitle{Shadow of the Future: Long-Term Thinking}
\begin{columns}
    \begin{column}[T]{0.5\linewidth}
      \begin{itemize}
        \item \textbf{Introduction}: Cooperation is sustainable only when the ``shadow of the future'' looms large. If interactions repeat indefinitely, the long-term reputational loss from cheating outweighs the short-term gain.
        \item \textbf{Situation}: Family businesses, small-town communities, subscription models, diplomatic relations, supplier relationships, wherever continuity of interaction creates accountability.
      \end{itemize}
    \end{column}
    \begin{column}[T]{0.5\linewidth}
      \begin{itemize}
        \item \textbf{Steps}:
        \begin{itemize}
          \item Increase the probability and salience of future interactions
          \item Reduce the discount rate on future payoffs
          \item Establish visible accountability systems
          \item Build reputational stakes into short-term interactions
          \item Create switching costs that make exit from the relationship costly
        \end{itemize}\item \textbf{Examples}: Not cheating your local vendor because you see them daily; subscription loyalty programs; diplomatic relations sustained by future needs; small-town honesty norms.
      \end{itemize}
    \end{column}
  \end{columns}
\end{frame}

%%%%%%%%%%%%%%%%%%%%%%%%%%%%%%%%%%%%%%%%%%%%%%%%%%%%%%%%%%%%%%%%%%%%%%%%%%%%%%%%%%
\begin{frame}[fragile]\frametitle{Signalling: Sending Messages Through Action}
\begin{columns}
    \begin{column}[T]{0.5\linewidth}
      \begin{itemize}
        \item \textbf{Introduction}: When one party holds private information, they must use observable actions to convey it credibly. A good signal must be both observable and difficult to fake, otherwise it conveys nothing.
        \item \textbf{Situation}: Job market signalling, brand advertising, corporate dividend policy, product warranties, personal reputation management, any information-asymmetric interaction.
      \end{itemize}
    \end{column}
    \begin{column}[T]{0.5\linewidth}
      \begin{itemize}
        \item \textbf{Steps}:
        \begin{itemize}
          \item Identify the hidden information you need to convey
          \item Find an action that is cheap for you (high type) but expensive for a low type to mimic
          \item Make the signal observable
          \item Maintain signal consistency over time
          \item Verify that recipients correctly interpret the signal
        \end{itemize}\item \textbf{Examples}: Wedding rings as availability signals; dividends signalling corporate financial health; peacock tails as evolutionary fitness signals; product warranties signalling quality.
      \end{itemize}
    \end{column}
  \end{columns}
\end{frame}

%%%%%%%%%%%%%%%%%%%%%%%%%%%%%%%%%%%%%%%%%%%%%%%%%%%%%%%%%%%%%%%%%%%%%%%%%%%%%%%%%%
\begin{frame}[fragile]\frametitle{Costly Signalling: Credibility Through Sacrifice}
\begin{columns}
    \begin{column}[T]{0.5\linewidth}
      \begin{itemize}
        \item \textbf{Introduction}: Talk is cheap; for a signal to be trusted in high-stakes games it must be too expensive for a ``faker'' to afford. The cost \textit{is} the proof.
        \item \textbf{Situation}: High-stakes relationship formation, luxury brand signalling, investor confidence building, gang initiation rites, religious fasting, any context where cheating would be profitable unless signalling is costly.
      \end{itemize}
    \end{column}
    \begin{column}[T]{0.5\linewidth}
      \begin{itemize}
        \item \textbf{Steps}:
        \begin{itemize}
          \item Identify what claim you need others to believe
          \item Find a costly action that only someone with the genuine claim would rationally take
          \item Ensure the cost is observable and verified
          \item Calibrate cost: high enough to deter fakers, not so high it deters honest signalers
          \item Repeat over time to build a track record
        \end{itemize}\item \textbf{Examples}: Expensive engagement rings; charitable donations by the wealthy; fraternity initiation costs; luxury goods as wealth signals; religious fasting as community commitment.
      \end{itemize}
    \end{column}
  \end{columns}
\end{frame}

%%%%%%%%%%%%%%%%%%%%%%%%%%%%%%%%%%%%%%%%%%%%%%%%%%%%%%%%%%%%%%%%%%%%%%%%%%%%%%%%%%
\begin{frame}[fragile]\frametitle{Rational Irrationality: Strategy Through Unpredictability}
\begin{columns}
    \begin{column}[T]{0.5\linewidth}
      \begin{itemize}
        \item \textbf{Introduction}: Being perfectly rational makes you predictable and exploitable. Feigning, or genuinely being, irrational alters the opponent's payoff matrix, forcing them to yield to avoid an unpredictable catastrophe.
        \item \textbf{Situation}: Geopolitical standoffs, extreme negotiation leverage, crowd control, deterring bullying, anywhere the opponent's certainty about your behaviour is a weapon against you.
      \end{itemize}
    \end{column}
    \begin{column}[T]{0.5\linewidth}
      \begin{itemize}
        \item \textbf{Steps}:
        \begin{itemize}
          \item Identify where your predictability is being exploited
          \item Introduce deliberate unpredictability or ``madman'' signals
          \item Ensure the irrationality is credible, genuine emotional investment helps
          \item Have a private off-switch you control even while appearing to lack one
          \item Return to rational behaviour once the leverage has worked
        \end{itemize}\item \textbf{Examples}: Nixon's ``Madman Theory'' in Vietnam; protester self-immolation as commitment device; doomsday devices in deterrence; road rage escalating beyond rational limits.
      \end{itemize}
    \end{column}
  \end{columns}
\end{frame}


% %%%%%%%%%%%%%%%%%%%%%%%%%%%%%%%%%%%%%%%%%%%%%%%%%%%%%%%%%%%%%%%%%%%%%%%%%%%%%%%%%%
% \begin{frame}[fragile]\frametitle{}
% \begin{center}
% {\Large Competitive Advantage and Market Strategy}
% \end{center}
% \end{frame}

%%%%%%%%%%%%%%%%%%%%%%%%%%%%%%%%%%%%%%%%%%%%%%%%%%%%%%%%%%%%%%%%%%%%%%%%%%%%%%%%%%
\begin{frame}[fragile]\frametitle{First-Mover Advantage: Being Early Pays}
\begin{columns}
    \begin{column}[T]{0.5\linewidth}
      \begin{itemize}
        \item \textbf{Introduction}: The first player to enter a market can lock in resources, customers, and standards. They set the rules that followers must accept. The brand can become the category name itself.
        \item \textbf{Situation}: New technology categories, emerging markets, standard-setting battles, patent races, geographic expansion into underpenetrated regions.
      \end{itemize}
    \end{column}
    \begin{column}[T]{0.5\linewidth}
      \begin{itemize}
        \item \textbf{Steps}:
        \begin{itemize}
          \item Identify the untapped market and timing window
          \item Move decisively before the window closes
          \item Use early position to lock in key resources (talent, patents, partners)
          \item Set the standard others must comply with
          \item Invest early in switching costs and network effects
        \end{itemize}\item \textbf{Examples}: Xerox defining ``photocopying''; Amazon in e-commerce; Coca-Cola's global brand position; Bitcoin as first cryptocurrency; squatters' rights laws.
      \end{itemize}
    \end{column}
  \end{columns}
\end{frame}

%%%%%%%%%%%%%%%%%%%%%%%%%%%%%%%%%%%%%%%%%%%%%%%%%%%%%%%%%%%%%%%%%%%%%%%%%%%%%%%%%%
\begin{frame}[fragile]\frametitle{Second-Mover Advantage: Learning from Mistakes}
\begin{columns}
    \begin{column}[T]{0.5\linewidth}
      \begin{itemize}
        \item \textbf{Introduction}: Being first is risky and expensive. Second movers free-ride on the pioneer's investments, learn from their errors, and adopt superior technology once the market is de-risked.
        \item \textbf{Situation}: Technology markets where innovation moves rapidly, emerging product categories after initial market education, regulatory environments that reward fast followers over pioneers.
      \end{itemize}
    \end{column}
    \begin{column}[T]{0.5\linewidth}
      \begin{itemize}
        \item \textbf{Steps}:
        \begin{itemize}
          \item Monitor pioneers carefully; let them educate the market
          \item Identify pioneer weaknesses and customer dissatisfactions
          \item Build a superior product on validated demand
          \item Strike at the moment of pioneer fatigue or overextension
          \item Use scale and capital efficiency to leapfrog
        \end{itemize}\item \textbf{Examples}: Google after Yahoo; Facebook after MySpace; Samsung after Apple; generic drugs after patent expiry; Zara replicating runway fashion rapidly.
      \end{itemize}
    \end{column}
  \end{columns}
\end{frame}

%%%%%%%%%%%%%%%%%%%%%%%%%%%%%%%%%%%%%%%%%%%%%%%%%%%%%%%%%%%%%%%%%%%%%%%%%%%%%%%%%%
\begin{frame}[fragile]\frametitle{Entry Deterrence: Protecting Market Territory}
\begin{columns}
    \begin{column}[T]{0.5\linewidth}
      \begin{itemize}
        \item \textbf{Introduction}: Incumbents can make it structurally unprofitable for new players to enter by over-investing in capacity, pricing aggressively, or controlling key complements, signalling ``don't bother.''
        \item \textbf{Situation}: Market leaders facing potential competition, platform companies with network effects, utilities and infrastructure businesses, industries with high capital requirements.
      \end{itemize}
    \end{column}
    \begin{column}[T]{0.5\linewidth}
      \begin{itemize}
        \item \textbf{Steps}:
        \begin{itemize}
          \item Signal excess capacity credibly to deter entry
          \item Use loyalty programs and switching costs to lock in customers
          \item Acquire critical input suppliers or distribution channels
          \item Use limit pricing to keep post-entry profits unattractive
          \item File strategic patents around core technology
        \end{itemize}\item \textbf{Examples}: Jio's free-data launch; frequent-flyer mile programmes; patent thickets in pharmaceuticals; exclusive distributor contracts; Amazon's low-margin saturation strategy.
      \end{itemize}
    \end{column}
  \end{columns}
\end{frame}

%%%%%%%%%%%%%%%%%%%%%%%%%%%%%%%%%%%%%%%%%%%%%%%%%%%%%%%%%%%%%%%%%%%%%%%%%%%%%%%%%%
\begin{frame}[fragile]\frametitle{Limit Pricing: Discouraging Competition Through Price}
\begin{columns}
    \begin{column}[T]{0.5\linewidth}
      \begin{itemize}
        \item \textbf{Introduction}: Setting prices lower than monopoly level, but above competitive cost, specifically to make market entry unprofitable for potential rivals without sacrificing all incumbent profits.
        \item \textbf{Situation}: Dominant firms in markets with moderate entry barriers, commodity markets with cost-advantage incumbents, digital platforms with significant scale economics.
      \end{itemize}
    \end{column}
    \begin{column}[T]{0.5\linewidth}
      \begin{itemize}
        \item \textbf{Steps}:
        \begin{itemize}
          \item Estimate the entrant's cost structure
          \item Set price just below the level that makes entry profitable for them
          \item Ensure price is above your own average cost
          \item Signal that price will stay low even post-entry
          \item Back the signal with visible cost advantages
        \end{itemize}\item \textbf{Examples}: OPEC pricing decisions; Microsoft Windows pricing in the 1990s; India's telecom data pricing wars; Amazon Basics undercutting third-party sellers.
      \end{itemize}
    \end{column}
  \end{columns}
\end{frame}

%%%%%%%%%%%%%%%%%%%%%%%%%%%%%%%%%%%%%%%%%%%%%%%%%%%%%%%%%%%%%%%%%%%%%%%%%%%%%%%%%%
\begin{frame}[fragile]\frametitle{Strategic Alliances: Temporary Partnerships}
\begin{columns}
    \begin{column}[T]{0.5\linewidth}
      \begin{itemize}
        \item \textbf{Introduction}: Cooperation between competitors, the ``frenemy'' relationship, to achieve specific goals like scale or access. Partners cooperate on the platform while competing on the product.
        \item \textbf{Situation}: Airline codeshare agreements, technology standards bodies, joint R\&D ventures, political coalitions, supply-chain partnerships with horizontal competitors.
      \end{itemize}
    \end{column}
    \begin{column}[T]{0.5\linewidth}
      \begin{itemize}
        \item \textbf{Steps}:
        \begin{itemize}
          \item Identify the specific shared objective that justifies cooperation
          \item Bound the alliance scope, cooperate only where interests align
          \item Establish clear governance and exit mechanisms
          \item Protect proprietary information carefully
          \item Plan for the alliance ending when objectives diverge
        \end{itemize}\item \textbf{Examples}: NATO; airline codeshare agreements; Spotify and Uber; Apple using Samsung screens; pharmaceutical joint ventures for R\&D.
      \end{itemize}
    \end{column}
  \end{columns}
\end{frame}

%%%%%%%%%%%%%%%%%%%%%%%%%%%%%%%%%%%%%%%%%%%%%%%%%%%%%%%%%%%%%%%%%%%%%%%%%%%%%%%%%%
\begin{frame}[fragile]\frametitle{Cheap Talk: Words Without Bite}
\begin{columns}
    \begin{column}[T]{0.5\linewidth}
      \begin{itemize}
        \item \textbf{Introduction}: Costless communication that does not directly affect payoffs. Because it is free, it is often misleading. But when interests align, cheap talk can successfully coordinate without commitment.
        \item \textbf{Situation}: Dating app profiles, political campaign promises, CEO letters, advertising claims, diplomatic statements, anywhere words are free and verification is costly.
      \end{itemize}
    \end{column}
    \begin{column}[T]{0.5\linewidth}
      \begin{itemize}
        \item \textbf{Steps}:
        \begin{itemize}
          \item Distinguish cheap talk from costly signalling
          \item Ask: does the speaker have incentives to lie? If yes, discount heavily
          \item Look for corroborating actions, not just words
          \item Use cheap talk to coordinate when interests align
          \item Design systems that make lying costly to convert cheap talk to credible talk
        \end{itemize}\item \textbf{Examples}: Street vendor's ``final offer''; political manifesto promises; CEO shareholder letters; football pre-match trash talk.
      \end{itemize}
    \end{column}
  \end{columns}
\end{frame}

% %%%%%%%%%%%%%%%%%%%%%%%%%%%%%%%%%%%%%%%%%%%%%%%%%%%%%%%%%%%%%%%%%%%%%%%%%%%%%%%%%%
% \begin{frame}[fragile]\frametitle{}
% \begin{center}
% {\Large Information and Uncertainty}
% \end{center}
% \end{frame}


%%%%%%%%%%%%%%%%%%%%%%%%%%%%%%%%%%%%%%%%%%%%%%%%%%%%%%%%%%%%%%%%%%%%%%%%%%%%%%%%%%
\begin{frame}[fragile]\frametitle{Perfect vs.\ Imperfect Information}
\begin{columns}
    \begin{column}[T]{0.5\linewidth}
      \begin{itemize}
        \item \textbf{Introduction}: Perfect information means all players see all moves and the full game state (chess). Imperfect information means some moves are hidden (poker). This single distinction transforms strategy from calculation to psychology.
        \item \textbf{Situation}: Distinguishing which type of game you are in determines whether to invest in deeper analysis vs.\ psychological profiling and information management.
      \end{itemize}
    \end{column}
    \begin{column}[T]{0.5\linewidth}
      \begin{itemize}
        \item \textbf{Steps}:
        \begin{itemize}
          \item Determine what information is observable vs.\ hidden
          \item In perfect information games, use backward induction and deep calculation
          \item In imperfect information games, model opponent types and beliefs
          \item Invest in revealing opponent information through screening
          \item Conceal your own information strategically
        \end{itemize}\item \textbf{Examples}: Chess (perfect) vs.\ Poker (imperfect); transparent vs.\ sealed-bid auctions; open parliamentary debates vs.\ backroom diplomacy.
      \end{itemize}
    \end{column}
  \end{columns}
\end{frame}

%%%%%%%%%%%%%%%%%%%%%%%%%%%%%%%%%%%%%%%%%%%%%%%%%%%%%%%%%%%%%%%%%%%%%%%%%%%%%%%%%%
\begin{frame}[fragile]\frametitle{Common Knowledge: Synchronised Understanding}
\begin{columns}
    \begin{column}[T]{0.5\linewidth}
      \begin{itemize}
        \item \textbf{Introduction}: Mutual knowledge means everyone knows X. Common knowledge means everyone knows that everyone knows X, to infinite recursion. This transition from private belief to shared awareness enables coordinated mass action.
        \item \textbf{Situation}: Social movements, bank runs, market bubbles, cultural paradigm shifts, anywhere private individual beliefs must become public to trigger collective behaviour change.
      \end{itemize}
    \end{column}
    \begin{column}[T]{0.5\linewidth}
      \begin{itemize}
        \item \textbf{Steps}:
        \begin{itemize}
          \item Identify information that is privately held but widely shared
          \item Find a public event that makes the private knowledge common
          \item Anticipate the coordinated behaviour change that follows
          \item Use or create common knowledge to coordinate allies
          \item Suppress common knowledge formation to prevent unwanted coordination
        \end{itemize}\item \textbf{Examples}: Arab Spring via social media; bank runs on shared panic; \#MeToo making isolated experiences collective; market bubbles bursting on shared doubt.
      \end{itemize}
    \end{column}
  \end{columns}
\end{frame}

%%%%%%%%%%%%%%%%%%%%%%%%%%%%%%%%%%%%%%%%%%%%%%%%%%%%%%%%%%%%%%%%%%%%%%%%%%%%%%%%%%
\begin{frame}[fragile]\frametitle{Information Asymmetry: The Knowledge Gap}
\begin{columns}
    \begin{column}[T]{0.5\linewidth}
      \begin{itemize}
        \item \textbf{Introduction}: When one party has more or better information than the other, the informed party can exploit the gap. This imbalance causes market failures as uninformed parties fear being cheated and exit.
        \item \textbf{Situation}: Doctor-patient, mechanic-driver, real estate agent-buyer, insurer-insured, employer-employee relationships, any principal-agent context with unequal information.
      \end{itemize}
    \end{column}
    \begin{column}[T]{0.5\linewidth}
      \begin{itemize}
        \item \textbf{Steps}:
        \begin{itemize}
          \item Identify who holds the information advantage
          \item If uninformed, seek signals, warranties, or third-party verification
          \item If informed, decide whether to disclose (building trust) or withhold (exploiting gap)
          \item Design mechanisms that align the informed party's incentives with disclosure
          \item Reduce asymmetry through transparency norms or regulation
        \end{itemize}\item \textbf{Examples}: Mechanics exploiting car-ignorant customers; insider trading; real estate agents prioritising commission over client interest.
      \end{itemize}
    \end{column}
  \end{columns}
\end{frame}

%%%%%%%%%%%%%%%%%%%%%%%%%%%%%%%%%%%%%%%%%%%%%%%%%%%%%%%%%%%%%%%%%%%%%%%%%%%%%%%%%%
\begin{frame}[fragile]\frametitle{Adverse Selection: Hidden Types Cause Problems}
\begin{columns}
    \begin{column}[T]{0.5\linewidth}
      \begin{itemize}
        \item \textbf{Introduction}: When quality is hidden, bad options (lemons) drive out good ones. Markets collapse because high-quality types leave rather than subsidise low-quality types, leaving only the worst in the pool.
        \item \textbf{Situation}: Insurance markets, used-car markets, online dating, volunteer armies, all-you-can-eat pricing, any market where participants self-select based on private information.
      \end{itemize}
    \end{column}
    \begin{column}[T]{0.5\linewidth}
      \begin{itemize}
        \item \textbf{Steps}:
        \begin{itemize}
          \item Identify the hidden characteristic causing selection
          \item Design screening mechanisms to sort types
          \item Use signalling to allow high types to distinguish themselves
          \item Price discriminate based on revealed preference
          \item Introduce mandatory pooling (e.g., compulsory insurance) to prevent market unravelling
        \end{itemize}\item \textbf{Examples}: Akerlof's used-car ``lemons'' market; health insurance death spirals; all-you-can-eat buffets attracting the most voracious eaters.
      \end{itemize}
    \end{column}
  \end{columns}
\end{frame}

%%%%%%%%%%%%%%%%%%%%%%%%%%%%%%%%%%%%%%%%%%%%%%%%%%%%%%%%%%%%%%%%%%%%%%%%%%%%%%%%%%
\begin{frame}[fragile]\frametitle{Moral Hazard: Risk Without Responsibility}
\begin{columns}
    \begin{column}[T]{0.5\linewidth}
      \begin{itemize}
        \item \textbf{Introduction}: When a person is insulated from the consequences of their actions, they behave more recklessly than if they bore the full risk. Shielding from downside destroys incentives for care.
        \item \textbf{Situation}: ``Too Big to Fail'' banks, comprehensively insured drivers, rent-controlled tenants, CEOs with golden parachutes, any principal-agent setting where consequences are borne by a third party.
      \end{itemize}
    \end{column}
    \begin{column}[T]{0.5\linewidth}
      \begin{itemize}
        \item \textbf{Steps}:
        \begin{itemize}
          \item Identify who bears the cost of risky behaviour
          \item Introduce co-payments, deductibles, or clawbacks to re-align incentives
          \item Increase monitoring to observe hidden actions
          \item Design contracts with performance-based compensation
          \item Use reputational accountability to supplement formal incentives
        \end{itemize}\item \textbf{Examples}: Banks gambling with taxpayer-backed deposits; drivers with comprehensive coverage taking more risks; CEOs with severance packages taking reckless bets.
      \end{itemize}
    \end{column}
  \end{columns}
\end{frame}

%%%%%%%%%%%%%%%%%%%%%%%%%%%%%%%%%%%%%%%%%%%%%%%%%%%%%%%%%%%%%%%%%%%%%%%%%%%%%%%%%%
\begin{frame}[fragile]\frametitle{Screening: Testing to Reveal Type}
\begin{columns}
    \begin{column}[T]{0.5\linewidth}
      \begin{itemize}
        \item \textbf{Introduction}: The uninformed party sets up a test that induces the informed party to reveal hidden information through self-selection. The design makes lying costly, so truth-telling becomes rational.
        \item \textbf{Situation}: Job interviews, insurance underwriting, university admissions, credit scoring, deductible-based insurance pricing, wherever the uninformed party needs to sort by hidden quality.
      \end{itemize}
    \end{column}
    \begin{column}[T]{0.5\linewidth}
      \begin{itemize}
        \item \textbf{Steps}:
        \begin{itemize}
          \item Identify the hidden type that matters (health, ability, commitment)
          \item Design a menu of options where each type prefers the option intended for them
          \item Ensure the ``wrong'' choice is costly enough to deter mimicry
          \item Observe self-selection and update beliefs
          \item Refine the screen as participants learn to game it
        \end{itemize}\item \textbf{Examples}: Insurance deductibles revealing health risk; Saturday-night stay airline pricing separating business from leisure; coding tests in job interviews.
      \end{itemize}
    \end{column}
  \end{columns}
\end{frame}

% %%%%%%%%%%%%%%%%%%%%%%%%%%%%%%%%%%%%%%%%%%%%%%%%%%%%%%%%%%%%%%%%%%%%%%%%%%%%%%%%%%
% \begin{frame}[fragile]\frametitle{}
% \begin{center}
% {\Large Negotiation and Bargaining}
% \end{center}
% \end{frame}



%%%%%%%%%%%%%%%%%%%%%%%%%%%%%%%%%%%%%%%%%%%%%%%%%%%%%%%%%%%%%%%%%%%%%%%%%%%%%%%%%%
\begin{frame}[fragile]\frametitle{Bargaining Power: Influencing Expectations}
\begin{columns}
    \begin{column}[T]{0.5\linewidth}
      \begin{itemize}
        \item \textbf{Introduction}: Your power in any negotiation is determined not by your loudness but by your BATNA, Best Alternative to a Negotiated Agreement. The party who can most credibly walk away controls the table.
        \item \textbf{Situation}: Salary negotiations, merger talks, hostage negotiations, real estate deals, supplier contracts, union collective bargaining, all negotiations where alternatives determine leverage.
      \end{itemize}
    \end{column}
    \begin{column}[T]{0.5\linewidth}
      \begin{itemize}
        \item \textbf{Steps}:
        \begin{itemize}
          \item Identify and actively improve your BATNA before negotiating
          \item Try to understand and weaken the other party's BATNA
          \item Never reveal a weak BATNA
          \item Signal your BATNA credibly through actions, not just words
          \item Anchor the negotiation range early with your opening offer
        \end{itemize}\item \textbf{Examples}: Having another job offer in salary negotiations; water vendor in a desert vs.\ one by a river; buyers in a buyer's market; monopolists with no alternatives facing them.
      \end{itemize}
    \end{column}
  \end{columns}
\end{frame}

%%%%%%%%%%%%%%%%%%%%%%%%%%%%%%%%%%%%%%%%%%%%%%%%%%%%%%%%%%%%%%%%%%%%%%%%%%%%%%%%%%
\begin{frame}[fragile]\frametitle{Nash Bargaining Solution: Fair Division}
\begin{columns}
    \begin{column}[T]{0.5\linewidth}
      \begin{itemize}
        \item \textbf{Introduction}: A mathematical approach to finding a ``fair'' bargaining outcome satisfying efficiency (nothing wasted), symmetry (equal bargainers get equal shares), and individual rationality. It creates a focal point for agreement.
        \item \textbf{Situation}: M\&A pricing, climate burden-sharing agreements, royalty splits, plea bargaining, joint-venture profit allocation, wherever a principled division of surplus is needed.
      \end{itemize}
    \end{column}
    \begin{column}[T]{0.5\linewidth}
      \begin{itemize}
        \item \textbf{Steps}:
        \begin{itemize}
          \item Identify each party's outside option (disagreement payoff)
          \item Calculate the total surplus from cooperation
          \item Find the split that maximises the product of surplus gains for both parties
          \item Check for symmetry adjustments based on relative patience or risk aversion
          \item Use the solution as a benchmark to evaluate proposed deals
        \end{itemize}\item \textbf{Examples}: M\&A deal pricing relative to stand-alone values; climate agreement burden-sharing; management-union wage settlements avoiding strikes.
      \end{itemize}
    \end{column}
  \end{columns}
\end{frame}

%%%%%%%%%%%%%%%%%%%%%%%%%%%%%%%%%%%%%%%%%%%%%%%%%%%%%%%%%%%%%%%%%%%%%%%%%%%%%%%%%%
\begin{frame}[fragile]\frametitle{Rubinstein Bargaining Model: Patience Wins}
\begin{columns}
    \begin{column}[T]{0.5\linewidth}
      \begin{itemize}
        \item \textbf{Introduction}: When two parties alternate offers over time, each delay costs both sides as the pie shrinks. But the more patient party captures a larger share. With rational players, the first offer is accepted immediately, delay signals irrational commitment or private information.
        \item \textbf{Situation}: Real estate negotiations with carrying costs, labour strikes, divorce proceedings, acquisition talks with dwindling startup runway, international trade negotiations.
      \end{itemize}
    \end{column}
    \begin{column}[T]{0.5\linewidth}
      \begin{itemize}
        \item \textbf{Steps}:
        \begin{itemize}
          \item Assess your discount rate, how much does delay cost you?
          \item Assess the opponent's discount rate, do they have burning runway?
          \item Make your first offer close to the equilibrium solution (it should be accepted)
          \item Signal patience credibly
          \item Recognise that real delays indicate incomplete information or irrational commitment, adjust accordingly
        \end{itemize}\item \textbf{Examples}: Car price negotiation as both parties' time costs mount; union strikes where both sides bleed daily; startup acquisition as funding runs out.
      \end{itemize}
    \end{column}
  \end{columns}
\end{frame}

%%%%%%%%%%%%%%%%%%%%%%%%%%%%%%%%%%%%%%%%%%%%%%%%%%%%%%%%%%%%%%%%%%%%%%%%%%%%%%%%%%
\begin{frame}[fragile]\frametitle{Rent-Seeking: Gaining Without Creating}
\begin{columns}
    \begin{column}[T]{0.5\linewidth}
      \begin{itemize}
        \item \textbf{Introduction}: Spending resources to increase one's share of existing wealth without creating new wealth. It is economically wasteful, consumes talent in zero-sum activities, and corrupts strategic incentives.
        \item \textbf{Situation}: Lobbying for regulations, patent trolling, occupational licensing cartels, corruption, feudalistic rent extraction, any activity focused on redistribution rather than value creation.
      \end{itemize}
    \end{column}
    \begin{column}[T]{0.5\linewidth}
      \begin{itemize}
        \item \textbf{Steps}:
        \begin{itemize}
          \item Distinguish between value-creating and rent-seeking activities in your strategy
          \item Calculate the full social cost of lobbying vs.\ innovating
          \item Identify rent-seeking by competitors and design regulation to block it
          \item Avoid engaging in rent-seeking even when profitable, it invites retaliation and regulatory scrutiny
          \item Advocate for institutional rules that reward production over position
        \end{itemize}\item \textbf{Examples}: India's ``Licence Raj'' lobbying; patent trolls; occupational licensing cartels; corruption and bribery as market entry barriers.
      \end{itemize}
    \end{column}
  \end{columns}
\end{frame}

% %%%%%%%%%%%%%%%%%%%%%%%%%%%%%%%%%%%%%%%%%%%%%%%%%%%%%%%%%%%%%%%%%%%%%%%%%%%%%%%%%%
% \begin{frame}[fragile]\frametitle{}
% \begin{center}
% {\Large Market Design and Mechanisms}
% \end{center}
% \end{frame}


%%%%%%%%%%%%%%%%%%%%%%%%%%%%%%%%%%%%%%%%%%%%%%%%%%%%%%%%%%%%%%%%%%%%%%%%%%%%%%%%%%
\begin{frame}[fragile]\frametitle{Mechanism Design: Engineering the Rules}
\begin{columns}
    \begin{column}[T]{0.5\linewidth}
      \begin{itemize}
        \item \textbf{Introduction}: Reverse game theory, instead of predicting what players do given rules, we ask what rules to design to elicit desired behaviour. It is the architecture of incentives, making truth-telling the best strategy.
        \item \textbf{Situation}: School choice systems, spectrum auctions, kidney exchange, tax code design, medical residency matching, any context where the designer controls the rules but not the players.
      \end{itemize}
    \end{column}
    \begin{column}[T]{0.5\linewidth}
      \begin{itemize}
        \item \textbf{Steps}:
        \begin{itemize}
          \item Define the social objective (efficiency, fairness, revenue)
          \item Identify the private information players hold
          \item Design a mechanism where revealing true information maximises individual payoffs (incentive-compatibility)
          \item Ensure participation is individually rational
          \item Test for manipulation and iterate
        \end{itemize}\item \textbf{Examples}: NRMP medical residency matching; FCC spectrum auctions; school choice algorithms; kidney exchange programmes; Vickrey second-price auctions.
      \end{itemize}
    \end{column}
  \end{columns}
\end{frame}

%%%%%%%%%%%%%%%%%%%%%%%%%%%%%%%%%%%%%%%%%%%%%%%%%%%%%%%%%%%%%%%%%%%%%%%%%%%%%%%%%%
\begin{frame}[fragile]\frametitle{Auction Theory: Strategic Bidding}
\begin{columns}
    \begin{column}[T]{0.5\linewidth}
      \begin{itemize}
        \item \textbf{Introduction}: Different auction formats (English, Dutch, sealed-bid, Vickrey) elicit different bidder behaviours and revenue outcomes. The winner's curse, overpaying because winning means everyone else bid lower, is the central hazard.
        \item \textbf{Situation}: IPL player auctions, government spectrum auctions, estate sales, eBay bidding, Google Ads, procurement auctions, privatisation of state assets.
      \end{itemize}
    \end{column}
    \begin{column}[T]{0.5\linewidth}
      \begin{itemize}
        \item \textbf{Steps}:
        \begin{itemize}
          \item Identify the auction format and its incentive structure
          \item Estimate your true private value, resist the winner's curse by bidding below your estimate in common-value auctions
          \item Observe competitor bidding patterns
          \item Use the Revenue Equivalence Theorem to evaluate format choices
          \item Bid truthfully in second-price (Vickrey) auctions, it is the dominant strategy
        \end{itemize}\item \textbf{Examples}: IPL team auctions for player talent; Google AdSense second-price mechanism; eBay online bidding; government bond issuance.
      \end{itemize}
    \end{column}
  \end{columns}
\end{frame}

%%%%%%%%%%%%%%%%%%%%%%%%%%%%%%%%%%%%%%%%%%%%%%%%%%%%%%%%%%%%%%%%%%%%%%%%%%%%%%%%%%
\begin{frame}[fragile]\frametitle{Matching Markets: Stable Allocations}
\begin{columns}
    \begin{column}[T]{0.5\linewidth}
      \begin{itemize}
        \item \textbf{Introduction}: Some markets cannot use prices to clear allocation (kidney donation, medical residency). Matching algorithms create stable outcomes where no unmatched pair would both prefer each other, preventing market unravelling and black markets.
        \item \textbf{Situation}: Medical residency assignment, school choice, kidney exchange networks, entry-level labour markets, organ donation registries, college admissions.
      \end{itemize}
    \end{column}
    \begin{column}[T]{0.5\linewidth}
      \begin{itemize}
        \item \textbf{Steps}:
        \begin{itemize}
          \item Verify that prices cannot or should not clear the market
          \item Gather complete preference lists from both sides
          \item Apply the Gale-Shapley deferred-acceptance algorithm
          \item Check that the resulting match is stable (no blocking pairs)
          \item Design incentives so participants reveal true preferences
        \end{itemize}\item \textbf{Examples}: NRMP medical residency (Nobel Prize-winning application); school choice in NYC and Boston; kidney exchange chains; entry-level lawyer placement.
      \end{itemize}
    \end{column}
  \end{columns}
\end{frame}

%%%%%%%%%%%%%%%%%%%%%%%%%%%%%%%%%%%%%%%%%%%%%%%%%%%%%%%%%%%%%%%%%%%%%%%%%%%%%%%%%%
\begin{frame}[fragile]\frametitle{Network Effects: Platforms and Lock-In}
\begin{columns}
    \begin{column}[T]{0.5\linewidth}
      \begin{itemize}
        \item \textbf{Introduction}: The value of a product increases with the number of its users. This creates winner-take-most dynamics where leading platforms grow larger while challengers die, not because the product is better but because the network is bigger.
        \item \textbf{Situation}: Social media, payment networks, operating systems, languages, stock exchanges, communication platforms, any product whose value scales with user count.
      \end{itemize}
    \end{column}
    \begin{column}[T]{0.5\linewidth}
      \begin{itemize}
        \item \textbf{Steps}:
        \begin{itemize}
          \item Measure same-side and cross-side network effects
          \item Subsidise early adoption to reach critical mass
          \item Design for interoperability or, strategically, against it
          \item Use data network effects to reinforce technical advantages
          \item Identify multi-homing costs, high multi-homing weakens lock-in
        \end{itemize}\item \textbf{Examples}: WhatsApp's dominance despite imperfect product; Visa/Mastercard duopoly; Windows OS market; English as global business language.
      \end{itemize}
    \end{column}
  \end{columns}
\end{frame}

%%%%%%%%%%%%%%%%%%%%%%%%%%%%%%%%%%%%%%%%%%%%%%%%%%%%%%%%%%%%%%%%%%%%%%%%%%%%%%%%%%
\begin{frame}[fragile]\frametitle{Two-Sided Markets: Balancing Multiple Sides}
\begin{columns}
    \begin{column}[T]{0.5\linewidth}
      \begin{itemize}
        \item \textbf{Introduction}: Platforms facilitating interaction between two distinct groups face a chicken-and-egg bootstrapping problem. The strategy requires subsidising one side to attract the other, then extracting value from whichever side has inelastic demand.
        \item \textbf{Situation}: Ride-sharing, dating apps, gaming consoles, newspapers, credit cards, shopping malls, any platform requiring both supply and demand side to create value.
      \end{itemize}
    \end{column}
    \begin{column}[T]{0.5\linewidth}
      \begin{itemize}
        \item \textbf{Steps}:
        \begin{itemize}
          \item Identify which side is the ``hard side'' to attract
          \item Subsidise the hard side heavily to build critical mass
          \item Charge the side with the most to gain from the other's presence
          \item Manage cross-side network effects, more of one side should benefit the other
          \item Prevent disintermediation, keep both sides transacting on the platform
        \end{itemize}\item \textbf{Examples}: Uber subsidising drivers and riders; gaming consoles subsidising hardware to attract developers; newspapers giving free content to attract advertisers.
      \end{itemize}
    \end{column}
  \end{columns}
\end{frame}


% %%%%%%%%%%%%%%%%%%%%%%%%%%%%%%%%%%%%%%%%%%%%%%%%%%%%%%%%%%%%%%%%%%%%%%%%%%%%%%%%%%
% \begin{frame}[fragile]\frametitle{}
% \begin{center}
% {\Large Collective Challenges and Social Dilemmas}
% \end{center}
% \end{frame}


%%%%%%%%%%%%%%%%%%%%%%%%%%%%%%%%%%%%%%%%%%%%%%%%%%%%%%%%%%%%%%%%%%%%%%%%%%%%%%%%%%
\begin{frame}[fragile]\frametitle{Tragedy of the Commons: Overuse Destroys Resources}
\begin{columns}
    \begin{column}[T]{0.5\linewidth}
      \begin{itemize}
        \item \textbf{Introduction}: When a shared resource is unregulated, individuals consume it to exhaustion, the benefit of consumption is private, but the cost of depletion is borne collectively. Rational individuals destroy the commons.
        \item \textbf{Situation}: Overfishing, groundwater depletion, traffic congestion, air pollution, office kitchen misuse, climate change, any shared resource without governance.
      \end{itemize}
    \end{column}
    \begin{column}[T]{0.5\linewidth}
      \begin{itemize}
        \item \textbf{Steps}:
        \begin{itemize}
          \item Make the resource excludable (property rights or licensing)
          \item Introduce quotas or Pigouvian taxes on usage
          \item Build community norms and social monitoring
          \item Apply Ostrom's principles: local rules, graduated sanctions, collective governance
          \item Create long-term stakes in the resource's health
        \end{itemize}\item \textbf{Examples}: Delhi's smog from uncurbed driving; global overfishing; groundwater depletion in agricultural regions; climate change as a global commons failure.
      \end{itemize}
    \end{column}
  \end{columns}
\end{frame}

%%%%%%%%%%%%%%%%%%%%%%%%%%%%%%%%%%%%%%%%%%%%%%%%%%%%%%%%%%%%%%%%%%%%%%%%%%%%%%%%%%
\begin{frame}[fragile]\frametitle{Voting Paradoxes: When Democracy Fails}
\begin{columns}
    \begin{column}[T]{0.5\linewidth}
      \begin{itemize}
        \item \textbf{Introduction}: Arrow's Impossibility Theorem proves that no rank-order voting system can perfectly satisfy all criteria of fairness simultaneously. Every voting system can be gamed or produce inconsistent collective choices.
        \item \textbf{Situation}: Political elections with multiple candidates, boardroom decisions, committee voting, Oscar nominations, budget prioritisation, any collective choice aggregating individual preferences.
      \end{itemize}
    \end{column}
    \begin{column}[T]{0.5\linewidth}
      \begin{itemize}
        \item \textbf{Steps}:
        \begin{itemize}
          \item Understand the voting mechanism and its incentives
          \item Identify potential spoiler effects and strategic voting incentives
          \item Consider alternative mechanisms (ranked-choice, approval voting)
          \item Account for agenda-setting power, who controls the order of votes
          \item Design voting systems to minimise strategic manipulation given the context
        \end{itemize}\item \textbf{Examples}: Spoiler candidates splitting the vote; gerrymandering; strategic ``lesser-evil'' voting; Condorcet cycles in committee decisions.
      \end{itemize}
    \end{column}
  \end{columns}
\end{frame}

%%%%%%%%%%%%%%%%%%%%%%%%%%%%%%%%%%%%%%%%%%%%%%%%%%%%%%%%%%%%%%%%%%%%%%%%%%%%%%%%%%
\begin{frame}[fragile]\frametitle{Volunteer's Dilemma: Who Acts First?}
\begin{columns}
    \begin{column}[T]{0.5\linewidth}
      \begin{itemize}
        \item \textbf{Introduction}: When a group needs one person to incur a cost for everyone's benefit, each member waits for another to volunteer. The larger the group, the more severe the problem, resulting in paralysis even when action is clearly needed.
        \item \textbf{Situation}: Emergency bystander responses, whistleblowing, group project initiation, social movements, bug reporting, standing ovations, reporting organisational wrongdoing.
      \end{itemize}
    \end{column}
    \begin{column}[T]{0.5\linewidth}
      \begin{itemize}
        \item \textbf{Steps}:
        \begin{itemize}
          \item Assign explicit responsibility to a specific person, diffused responsibility equals paralysis
          \item Reward first-movers visibly and publicly
          \item Reduce the cost of volunteering (anonymity, indemnification)
          \item Break large groups into smaller accountable units
          \item Build a culture where volunteering is the expected norm
        \end{itemize}\item \textbf{Examples}: Bystander effect in the Kitty Genovese case; whistleblowers in corporations; group projects where no one starts; first protesters in social movements.
      \end{itemize}
    \end{column}
  \end{columns}
\end{frame}

%%%%%%%%%%%%%%%%%%%%%%%%%%%%%%%%%%%%%%%%%%%%%%%%%%%%%%%%%%%%%%%%%%%%%%%%%%%%%%%%%%
\begin{frame}[fragile]\frametitle{Focal Points (Schelling Points): Coordinating Without Communication}
\begin{columns}
    \begin{column}[T]{0.5\linewidth}
      \begin{itemize}
        \item \textbf{Introduction}: When multiple equilibria exist and communication is impossible, people converge on solutions that stand out through prominence, symmetry, or cultural salience. Self-fulfilling expectations create convergence without explicit agreement.
        \item \textbf{Situation}: Emergency rendezvous, social conventions, negotiation anchors, ceasefire boundary selection, standard pricing points, any coordination without the ability to communicate.
      \end{itemize}
    \end{column}
    \begin{column}[T]{0.5\linewidth}
      \begin{itemize}
        \item \textbf{Steps}:
        \begin{itemize}
          \item Identify which options stand out as ``obvious'' in context
          \item Use round numbers and symmetric solutions as coordination anchors
          \item In design, make your desired outcome the salient focal point
          \item Use cultural and historical reference points to guide coordination
          \item Recognise that your opponents will also search for the same focal point
        \end{itemize}\item \textbf{Examples}: Grand Central Station at noon as a meeting point; round-number salary anchors (Rs. 10 lakh, not Rs. 9.83 lakh); ceasefire lines at rivers; splitting bills equally.
      \end{itemize}
    \end{column}
  \end{columns}
\end{frame}

%%%%%%%%%%%%%%%%%%%%%%%%%%%%%%%%%%%%%%%%%%%%%%%%%%%%%%%%%%%%%%%%%%%%%%%%%%%%%%%%%%
\begin{frame}[fragile]\frametitle{The Art of War and Diplomacy: Ancient Strategic Wisdom}
\begin{columns}
    \begin{column}[T]{0.5\linewidth}
      \begin{itemize}
        \item \textbf{Introduction}: Chanakya and Sun Tzu wrote millennia ago, yet CEOs quote them today. Because technology changes but human nature and conflict do not. The highest strategy is to win without fighting.
        \item \textbf{Situation}: Corporate competition, geopolitical diplomacy, organisational politics, competitive intelligence, brand wars, wherever perception management and asymmetric information can substitute for direct confrontation.
      \end{itemize}
    \end{column}
    \begin{column}[T]{0.5\linewidth}
      \begin{itemize}
        \item \textbf{Steps}:
        \begin{itemize}
          \item Know yourself and your enemy, assess strengths and weaknesses honestly
          \item Win without fighting: use perception, reputation, and deception before force
          \item Choose the battlefield, fight where you are strong, not where the enemy is
          \item Use spies and intelligence before commitment
          \item Strike at the opponent's strategy, not their army
        \end{itemize}\item \textbf{Examples}: Corporate hostile takeover via proxy; soft power (Bollywood/Hollywood); guerrilla marketing; peace treaties as strategic positioning; cyber warfare.
      \end{itemize}
    \end{column}
  \end{columns}
\end{frame}

% %%%%%%%%%%%%%%%%%%%%%%%%%%%%%%%%%%%%%%%%%%%%%%%%%%%%%%%%%%%%%%%%%%%%%%%%%%%%%%%%%%
% \begin{frame}[fragile]\frametitle{}
% \begin{center}
% {\Large Prospect Theory}
% \end{center}
% \end{frame}



%%%%%%%%%%%%%%%%%%%%%%%%%%%%%%%%%%%%%%%%%%%%%%%%%%%%%%%%%%%%%%%%%%%%%%%%%%%%%%%%%%
\begin{frame}[fragile]\frametitle{Prospect Theory: Loss Aversion Changes Everything}
\begin{columns}
    \begin{column}[T]{0.45\linewidth}
      \begin{itemize}
        \item \textbf{Introduction}: Kahneman and Tversky's theory shows humans evaluate outcomes relative to a reference point, not absolute wealth, and losses hurt roughly twice as much as equivalent gains feel good. This reshapes all strategic interaction.
        \item \textbf{Situation}: Investor behaviour, labour negotiations, political framing, consumer pricing, legal settlements, any context where framing around a reference point shifts decisions away from rational maximisation.
      \end{itemize}
    \end{column}
    \begin{column}[T]{0.55\linewidth}
      \begin{itemize}
        \item \textbf{Steps}:
        \begin{itemize}
          \item Identify the reference point of your counterpart
          \item Frame your offer as avoiding a loss rather than achieving a gain
          \item Recognise the endowment effect, people overvalue what they already possess
          \item Avoid triggering loss aversion by anchoring to a favourable reference point early
          \item Exploit loss aversion in competitors while designing your own decisions around rational expected value
        \end{itemize}\item \textbf{Examples}: Investors refusing to sell losing stocks; homeowners overpricing due to emotional attachment; unions rejecting wage cuts even to prevent layoffs; governments escalating losing wars to ``not lose face.''
      \end{itemize}
    \end{column}
  \end{columns}
\end{frame}

%%%%%%%%%%%%%%%%%%%%%%%%%%%%%%%%%%%%%%%%%%%%%%%%%%%%%%%%%%%%%%%%%%%%%%%%%%%%%%%%%%
\begin{frame}[fragile]\frametitle{}
\begin{center}
{\Large Conclusion and Next Steps}
\end{center}
\end{frame}


%%%%%%%%%%%%%%%%%%%%%%%%%%%%%%%%%%%%%%%%%%%%%%%%%%%%%%%%%%%%%%%%%%%%%%%%%%%%%%%%%%
\begin{frame}[fragile]\frametitle{The Infinite Game: Key Takeaways}
  \begin{columns}
    \begin{column}[T]{0.48\linewidth}
      \begin{block}{What We Learned}
        \begin{itemize}
          \item You are never playing solitaire, always chess
          \item Incentives trump intentions every time
          \item Information is the most valuable commodity
          \item Credibility must be earned through action
          \item Cooperation requires design, not hope
          \item Reputation compounds over decades
          \item The framing of a choice changes the choice
        \end{itemize}
      \end{block}
    \end{column}
    \begin{column}[T]{0.48\linewidth}
      \begin{block}{The Hierarchy of Strategic Wisdom}
        \begin{itemize}
          \item Know which game you are in
          \item Know who the other players are
          \item Know their incentives
          \item Know your BATNA
          \item Make credible commitments
          \item Build reputation relentlessly
          \item \textbf{Play to keep playing}
        \end{itemize}
      \end{block}
  \vspace{8pt}
  \begin{alertblock}{The Ultimate Lesson}
    In the \textit{Finite Game} of a football match, you play to win. In the \textit{Infinite Game} of life, you play to \textbf{keep playing}. Stay adaptable.
  \end{alertblock}	  
    \end{column}
  \end{columns}

\end{frame}

%%%%%%%%%%%%%%%%%%%%%%%%%%%%%%%%%%%%%%%%%%%%%%%%%%%%%%%%%%%%%%%%%%%%%%%%%%%%%%%%%%
\begin{frame}[fragile]\frametitle{Applying Game Theory: Your Action Plan}
  \begin{columns}
    \begin{column}[T]{0.48\linewidth}
      \begin{block}{This Week}
        \begin{itemize}
          \item Map one negotiation you are in as a game tree
          \item Identify your BATNA and improve it
          \item Find one Nash Equilibrium you are stuck in and ask: \textit{what changes the game?}
          \item Spot one dominant strategy and use it
        \end{itemize}
      \end{block}


    \end{column}
    \begin{column}[T]{0.48\linewidth}
      \begin{block}{This Month}
        \begin{itemize}
          \item Redesign one professional relationship using Tit-for-Tat
          \item Audit your reputation signals
          \item Find one rent-seeking activity and eliminate it
          \item Implement one pre-commitment device for a personal goal
        \end{itemize}
      \end{block}
    \end{column}
  \end{columns}
\end{frame}

%%%%%%%%%%%%%%%%%%%%%%%%%%%%%%%%%%%%%%%%%%%%%%%%%%%%%%%%%%%%%%%%%%%%%%%%%%%%%%%%%%
\begin{frame}[fragile]\frametitle{Applying Game Theory: Your Action Plan}
  \begin{columns}
    \begin{column}[T]{0.48\linewidth}
      \begin{block}{Ongoing Practice}
        \begin{itemize}
          \item Before every significant decision: \textit{who else is deciding, and what are their payoffs?}
          \item In every negotiation: \textit{what is the other party's reference point?}
          \item In every conflict: \textit{is this zero-sum or can I grow the pie?}
          \item In every commitment: \textit{is this threat or promise actually credible?}
        \end{itemize}
      \end{block}
    \end{column}
    \begin{column}[T]{0.48\linewidth}

      \begin{alertblock}{Remember}
        Game Theory is a map, not the territory. Real people are irrational, emotional, and surprising. Use these models as lenses, not algorithms.
      \end{alertblock}
    \end{column}
  \end{columns}
\end{frame}

%%%%%%%%%%%%%%%%%%%%%%%%%%%%%%%%%%%%%%%%%%%%%%%%%%%%%%%%%%%%%%%%%%%%%%%%%%%%%%%%%%
\begin{frame}[fragile]\frametitle{Essential Reading: Resources to Go Deeper}
  \begin{columns}
    \begin{column}[T]{0.48\linewidth}
      \begin{block}{Foundational Books}
        \begin{itemize}
          \item \textbf{Thinking Strategically}, Dixit \& Nalebuff (best accessible intro)
          \item \textbf{The Art of Strategy}, Dixit \& Nalebuff (applied game theory)
          \item \textbf{Games of Strategy}, Dixit, Skeath \& McAdams (textbook)
          \item \textbf{The Strategy of Conflict}, Thomas Schelling (focal points, brinkmanship)
          \item \textbf{Thinking, Fast and Slow}, Daniel Kahneman (Prospect Theory)
          \item \textbf{Getting to Yes}, Fisher \& Ury (applied negotiation)
        \end{itemize}
      \end{block}
    \end{column}
    \begin{column}[T]{0.48\linewidth}
      \begin{block}{Advanced \& Specialised}
        \begin{itemize}
          \item \textbf{Arthashastra}, Kautilya (Chanakya), ancient strategy
          \item \textbf{The Art of War}, Sun Tzu, conflict \& deception
          \item \textbf{Predictably Irrational}, Dan Ariely, behavioural economics
          \item \textbf{Co-opetition}, Brandenburger \& Nalebuff, business strategy
          \item \textbf{The Infinite Game}, Simon Sinek, finite vs.\ infinite perspective
          \item \textbf{Who Gets What and Why}, Alvin Roth, matching markets (Nobel)
        \end{itemize}
      \end{block}
      % \begin{block}{Online Courses}
        % \begin{itemize}
          % \item Yale / Coursera: \textit{Game Theory} (Ben Polak)
          % \item Stanford / Coursera: \textit{Game Theory} (Jackson, Shoham, Leyton-Brown)
          % \item MIT OpenCourseWare: 14.12 Economic Applications of Game Theory
        % \end{itemize}
      % \end{block}
    \end{column}
  \end{columns}
\end{frame}

%%%%%%%%%%%%%%%%%%%%%%%%%%%%%%%%%%%%%%%%%%%%%%%%%%%%%%%%%%%%%%%%%%%%%%%%%%%%%%%%%%
\begin{frame}[fragile]\frametitle{Closing Reflection}

\begin{block}{Thinking Strategically \ldots}
  \begin{itemize}
    \item ``The supreme art of war is to subdue the enemy 
      without fighting.'' - Sun Tzu, The Art of War
	\item ``In the real world, the game isn't broken.   The players are human.'' - Kahneman / Tversky, Prospect Theory Reflection
    \item Understand the incentives.
    \item Respect the players.
    \item May your equilibrium always be stable.
  \end{itemize}
\end{block}  
\end{frame}